% Mechanically extracted after \mainmatter from the frozen standalone master.

\markboth{Exposé XVI. Groups of Unipotent Rank Zero}
  {Exposé XVI. Groups of Unipotent Rank Zero}
% BEGIN inputs/XVI/tex/components/01_expose_XVI_section1_immersion_criterion.tex
% SGA 3, Exposé XVI, Loop-1 English reconstruction.
% Current scope: title and all of Section 1.
% Authority: Polo--Gille Expo16.pdf, local pages 1--10
% (combined-reader pages 939--948; printed pages 484--502).
% Local GPU OCR and the Reinhold English Markdown are comparison/locator
% witnesses only. The Polo--Gille page image controls text and formulas.

\setcounter{footnote}{0}
\SGARefTarget{sga3:xvi:chapter:xvi}\chapter*{Exposé XVI\\
Groups of Unipotent Rank Zero}
\addcontentsline{toc}{chapter}{Exposé XVI. Groups of Unipotent Rank Zero}
\markboth{EXPOSÉ XVI. GROUPS OF UNIPOTENT RANK ZERO}
  {EXPOSÉ XVI. GROUPS OF UNIPOTENT RANK ZERO}

\begin{center}
  \textit{by M. Raynaud}\textsuperscript{(*)}
\end{center}
\begingroup
\renewcommand{\thefootnote}{0}
\footnotetext{xy version of December 1, 2008.}
\renewcommand{\thefootnote}{*}
\footnotetext{Cf.\ the note on page 1 of Exposé XV.}
\endgroup

\SGARefTarget{sga3:xvi:section:1}\section*{1. An Immersion Criterion}
\addcontentsline{toc}{section}{1. An Immersion Criterion}

\SGARefTarget{sga3:xvi:subsection:1:1}%
\subsection*{1.1. Examples of Monomorphisms of Group Preschemes
That Are Not Immersions}
\addcontentsline{toc}{subsection}{\SGARefLink{sga3:xvi:subsection:1:1}{1.1}. Examples of Monomorphisms of
Group Preschemes That Are Not Immersions}

We shall construct a prescheme \(S\), two \(S\)-group preschemes
\(G\) and \(H\), and an \(S\)-group monomorphism
\[
  u:G\longrightarrow H
\]
that is not an immersion. The groups \(G\) and \(H\) will be of finite
presentation over \(S\), flat over \(S\), and \(G\) will even be
smooth over \(S\) (cf.\ Exposé VIII, the paragraph and note~\((*)\)
preceding \SGARefLink{sga3:viii:lemma:7:1}{7.1}; see also XVII, Appendix III, 4).

First take \(S\) to be the spectrum of a discrete valuation ring
\(A\) of unequal characteristic and residue characteristic \(2\).
Let \(t\) be the generic point of \(S\) and \(s\) the closed point.

\medskip
\noindent\textit{a)} Take for \(G\) the open subgroup of the constant
group
\[
  G'=(\mathbb Z/2\mathbb Z)_S
\]
obtained by removing the point of the closed fiber
\((\mathbb Z/2\mathbb Z)_s\) distinct from the identity element, and
take for \(H\) the group of multiplicative type
\((\mu_2)_S\) of square roots of unity (\SGARefLink{sga3:i:paragraph:4-4-4}{Exposé I, 4.4.4}). By the
choice of \(S\), \(H_t\) is isomorphic to
\((\mathbb Z/2\mathbb Z)_t\), whereas \(H_s\) is a radicial group.
There is an evident morphism \(G'\to H\), defined by the section
\((-1)\) of \(H\), hence a morphism \(u:G\to H\) which is a
monomorphism but not an immersion (otherwise it would necessarily be
an isomorphism; cf.\ \SGARefLink{sga3:viii:section:7}{Exposé VIII, 7}). Here \SGARefLink{sga3:viii:proposition:7:9}{Exposé VIII, 7.9} does not
apply, since \({}_{2}G=G\) is not finite over \(S\).

\medskip
\noindent\textit{b)} Let \(K\) denote the étale, nonseparated
\(S\)-group prescheme obtained from the identity group by
``doubling'' the unique point of the closed fiber. Take for \(H\) the
product group
\[
  (\mu_2)_S\times_S K.
\]
Let \(a\) (respectively \(b\)) be the unique section of \(\mu_2\)
(respectively \(K\)) over \(S\) distinct from the identity section,
and let \(h=(a,b)\) be the corresponding section of \(H\). The datum
of \(h\) defines an \(S\)-group morphism
\[
  u:G=(\mathbb Z/2\mathbb Z)_S\longrightarrow H.
\]
It is clear that \(u\) is a monomorphism and is not an immersion.
Here \SGARefLink{sga3:viii:proposition:7:9}{Exposé VIII, 7.9} does not apply, since \({}_{2}H=H\) is not
separated over \(S\).

\medskip
\noindent\SGAInlineRefTarget{sga3:xvi:subsection:1:1:item:c}\textit{c)} More interesting is the following example, in
which \(G\) is a smooth group with connected fibers (and even
\(G=(\Ga)_S\)); cf.\ \SGARefLink{sga3:viii:remark:7:10}{Exposé VIII, 7.10}.

Take for \(S\) the spectrum of a discrete valuation ring \(A\) of
equal characteristic \(2\), and let \(\pi\) be a uniformizer of
\(A\) and \(\mathfrak m\) the maximal ideal of \(A\).

Consider the additive group \((\Ga)_S\), the product group
\((\Ga\times\Ga)_S\) with ring \(A[X,Y]\), and let \(G_1\) be the
closed subgroup defined by
\begin{equation}
  X^2+X-\pi Y=0.
  \tag{1}\label{sga3:xvi:eq:1}
\end{equation}
The group \(G_1\) is smooth over \(S\); its generic fiber is
isomorphic to \((\Ga)_t\), and its closed fiber is isomorphic to the
product of \((\Ga)_s\) with \((\mathbb Z/2\mathbb Z)_s\). For the
\((\mathbb Z/2\mathbb Z)_s\)-factor one may take the subgroup \(N_s\)
of \((G_1)_s\) defined by
\begin{equation}
  Y=X^2+X=0.
  \tag{2}\label{sga3:xvi:eq:2}
\end{equation}

Let \(N'\) and \(N''\) be two group subschemes of \(G_1\), isomorphic
to \((\mathbb Z/2\mathbb Z)_S\), whose closed fibers coincide with
\(N_s\) and whose generic fibers are distinct. Thus \(N'\) and
\(N''\) are defined by two sections of \(G_1\) over \(S\), with
coordinates \((a',b')\) and \((a'',b'')\), respectively, such that
\begin{equation}
  \begin{cases}
    a'^2+a'-\pi b'=a''^2+a''-\pi b''=0,\\
    a'\text{ and }a''\equiv 1\pmod{\mathfrak m},
      \qquad a'\ne a'',\\
    b'\text{ and }b''\equiv 0\pmod{\mathfrak m}.
  \end{cases}
  \tag{3}\label{sga3:xvi:eq:3}
\end{equation}
(For example, one may take \((1,0)\) and
\((1+\pi^2,\pi+\pi^3)\).)

The quotient groups
\[
  G'=G_1/N'
  \qquad\text{and}\qquad
  G''=G_1/N''
\]
are representable (\SGARefLink{thm:sga3-V-4.1}{Exposé V, 4.1}). Let \(u\) be the canonical
morphism
\[
  u:G_1\longrightarrow G'\times_S G'',
\]
and let \(H\) be the group subscheme of \(G'\times_S G''\) equal to
the schematic closure in \(G'\times_S G''\) of
\(u_t((G_1)_t)\), so that \(u\) factors through \(H\). The kernel of
\(u\) is
\[
  K=N'\cap N'',
\]
whose generic fiber is the identity group and whose closed fiber is
equal to \(N_s\); in particular, \(K\) is not flat over \(S\). On the
generic fiber, \(u_t\) is therefore an isomorphism from
\((G_1)_t\) onto \(H_t\). On the other hand, \(H_s\) is not smooth.
Indeed, if \(H_s\) were smooth, then, since \(u_s\) is a morphism with
finite kernel and \((G_1)_s\) and \(H_s\) have the same dimension,
namely \(1\), \(u_s\) would be flat. Since \(G_1\) and \(H\) are flat,
\(u\) would then be flat, and consequently \(\Ker u\) would be flat
over \(S\), which it is not. It is clear that the restriction of
\(u\) to the identity component \(G\) of \(G_1\) is a monomorphism:
the fibers of \(K\cap G\) are identity groups, so \(K\cap G\) is the
identity group (Exposé VI\textsubscript{B}, \SGARefLink{lem:sga3-VIB-2.9}{2.9}). But it is not an
immersion, for otherwise it would be an open immersion and \(H_s\)
would be smooth. Here \SGARefLink{sga3:viii:proposition:7:9}{Exposé VIII, 7.9} does not apply, since
\({}_{2}G=G\) is not finite over \(S\).

For lovers of equations, in the preceding example one may take for
\(H\) the closed subgroup of \((\Ga\times_S\Ga)_S\) with ring
\(A[V,W]/(F)\), where
\[
  F=(a''b'-b''a')(a''^2V-a'^2W)-(a''V-a'W)^2
\]
(where \(a',a'',b',b''\) satisfy \eqref{sga3:xvi:eq:3}). For \(G\)
one takes the group \((\Ga)_S\) with ring \(A[T]\), and for the
morphism \(u:G\to H\) the \(S\)-morphism defined by
\[
  V\longmapsto a'(a'T+\pi T^2),
  \qquad
  W\longmapsto a''(a''T+\pi T^2).
\]

\medskip
\noindent\SGARefTarget{sga3:xvi:remark:1:2}%
\textbf{Remark 1.2.} --- The preceding construction is inspired by
the method of Koizumi--Shimura.\footnote{Editor's note:
S.~Koizumi and G.~Shimura, ``Specialization of abelian varieties,''
\emph{Sci.\ Papers Coll.\ Gen.\ Ed.\ Univ.\ Tokyo} 9 (1959),
187--211.} It does not apply when \(S\) is the spectrum of a ring of
unequal characteristic, because the points of finite order of \(G\)
are then ``too rigid.''

\SGARefTarget{sga3:xvi:subsection:1:3}%
\subsection*{1.3. Statement of the Immersion Criterion}
\addcontentsline{toc}{subsection}{1.3. Statement of the Immersion Criterion}

\medskip
\noindent\SGARefTarget{sga3:xvi:theorem:1:3}%
\textbf{Theorem 1.3.} --- \textit{Let \(S\) be a locally noetherian
prescheme, let \(G\) be an \(S\)-group prescheme, smooth and of finite
type over \(S\), possessing Cartan subgroups locally for the fpqc
topology (\SGARefLink{sga3:xv:definition:6-1}{Exposé XV, 6.1} and \SGARefLink{sga3:xv:theorem:7-3}{7.3}~(i)), let \(H\) be an \(S\)-group
prescheme locally of finite type, and let \(u:G\to H\) be an
\(S\)-group monomorphism. Then:}

\begin{enumerate}
\item[\textit{a)}] \textit{If \(G\) has connected fibers, then \(u\)
is an immersion if and only if, for every \(S\)-scheme \(S'\) that is
the spectrum of a complete discrete valuation ring with algebraically
closed residue field, and every Cartan subgroup \(C\) of \(G_{S'}\),
the restriction of \(u_{S'}\) to \(C\) is an immersion.}

\item[\textit{b)}] \textit{The morphism \(u\) is an immersion
(respectively a closed immersion) if and only if, for every \(S'\) as
above and every Cartan subgroup \(C\) of the identity component
\((G_{S'})^0\) of \(G_{S'}\), the restriction of \(u_{S'}\) to the
normalizer \(N\) of \(C\) in \(G_{S'}\) (\SGARefLink{sga3:xi:corollary:6:11}{Exposé XI, 6.11}) is an
immersion (respectively a closed immersion).}
\end{enumerate}

Before proving \SGARefLink{sga3:xvi:theorem:1:3}{1.3}, let us state some applications.

\medskip
\noindent\SGARefTarget{sga3:xvi:corollary:1:4}%
\textbf{Corollary 1.4.} --- \textit{Let \(S\) be a prescheme, let
\(G\) be an \(S\)-group prescheme, smooth and of finite presentation
over \(S\), of unipotent rank zero (\SGARefLink{sga3:xv:definition:6:1:ter}{Exposé XV, 6.1 ter}), possessing
maximal tori locally for the fpqc topology, let \(H\) be an
\(S\)-group prescheme, and let \(u:G\to H\) be an \(S\)-group
monomorphism. Suppose in addition either that \(H\) is of finite
presentation over \(S\), or that \(S\) is locally noetherian and
\(H\) is locally of finite type. Then:}

\begin{enumerate}
\item[\textit{a)}] \textit{If \(G\) has connected fibers, \(u\) is
an immersion.}

\item[\textit{b)}] \textit{If \(H\) is separated over \(S\), and if
for every \(S'\) over \(S\) and every maximal torus \(T\) of
\(G_{S'}\), the Weyl group}
\[
  W=
  \Norm_{G_{S'}}(T)\big/\Centr_{G_{S'}^0}(T)
\]
\textit{is representable (a condition always satisfied when \(G\)
has connected fibers; \SGARefLink{sga3:xv:proposition:7-1}{Exposé XV, 7.1}~(iv)) and finite over \(S\), then
\(u\) is a closed immersion.}
\end{enumerate}

\medskip
\noindent\SGARefTarget{sga3:xvi:corollary:1:5}%
\textbf{Corollary 1.5.} --- \textit{Let \(S\) be a prescheme, let
\(G\) be an \(S\)-group prescheme, smooth and of finite presentation
over \(S\), with connected fibers, let \(H\) be an \(S\)-group
prescheme, and let \(u:G\to H\) be an \(S\)-group monomorphism.
Suppose either that \(H\) is of finite presentation over \(S\), or
that \(S\) is locally noetherian and \(H\) is locally of finite type.
Then:}

\begin{enumerate}
\item[\textit{a)}] \textit{If \(G\) is reductive, that is, affine
over \(S\) with reductive fibers (cf.\ Exposé XIX), \(u\) is an
immersion, and a closed immersion if \(H\) is separated.}

\item[\textit{b)}] \textit{If \(G\) has affine, solvable, connected
fibers of unipotent rank zero, \(u\) is an immersion, and a closed
immersion if \(H\) is separated.}
\end{enumerate}

\medskip
\noindent\textit{Proof of \SGARefLink{sga3:xvi:corollary:1:4}{1.4} from \SGARefLink{sga3:xvi:theorem:1:3}{1.3}.}

\smallskip
\noindent\textit{Reduction to the noetherian case.}
If \(S\) is locally noetherian, the reduction is immediate, since the
properties to be proved are local on \(S\). In the second case \(G\)
and \(H\) are of finite presentation over \(S\). After shrinking
\(S\), we may suppose that \(S\) is affine. By Exposé
VI\textsubscript{B}, \S10, there exist a noetherian scheme \(S_0\), an
\(S_0\)-group prescheme \(G_0\), smooth over \(S_0\) (with connected
fibers in case~a)) and of finite type over \(S_0\), an
\(S_0\)-group prescheme \(H_0\) of finite type (separated in case~b)),
and an \(S_0\)-group morphism
\[
  u_0:G_0\longrightarrow H_0
\]
such that \(G,H,u\) are obtained from \(G_0,H_0,u_0\) by the base
extension \(S\to S_0\). The assertion that \(u\) is a monomorphism
means that \(\Ker u\) is the identity group; we may therefore suppose
that \(u_0\) is a monomorphism.

Since \(G\) has unipotent rank \(\rho_u=0\) and possesses maximal
tori locally for the fpqc topology, the abelian rank \(\rho_{ab}\) of
the fibers of \(G\) is locally constant (\SGARefLink{sga3:xv:corollary:8-18}{Exposé XV, 8.18}). But
\(\rho_u\) and \(\rho_{ab}\) are locally constructible functions
(Exposé XV, 6.3 bis). A standard argument (cf.\ EGA IV, 8.3.4) shows
that one may choose \(S_0\) and \(G_0\) so that the unipotent rank
(respectively the abelian rank) of the fibers of \(G_0\) is zero
(respectively locally constant). Then \(G_0\) possesses maximal tori
locally for the fpqc topology (\SGARefLink{sga3:xv:corollary:8-18}{Exposé XV, 8.18}), and the functor
\(\mathcal{TM}_{G_0}\) of maximal tori of \(G_0\) is representable by
an \(S_0\)-scheme \(X_0\) of finite type over \(S_0\) (\SGARefLink{sga3:xv:theorem:8-15}{Exposé XV,
8.15}). Let \(T_0\) be the ``universal'' maximal torus of
\((G_0)_{X_0}\), let \(X=X_0\times_{S_0}S\), and let
\(T=T_0\times_{S_0}S\), the universal maximal torus for \(G\). By
hypothesis, in case~b), the Weyl group relative to \(T\) is
representable and finite over \(X\). These two properties are
compatible with projective limits of preschemes (Exposé
VI\textsubscript{B}, 10.1~iii, and EGA IV, 8.10.5). We may therefore
choose \(S_0\) so that the Weyl group of \(T_0\) is finite over
\(X_0\). Under these conditions, to prove \SGARefLink{sga3:xvi:corollary:1:4}{1.4} we may replace
\(S,G,u,H\) by \(S_0,G_0,u_0,H_0\), and hence suppose that \(S\) is
noetherian.

Use the valuative criterion furnished by \SGARefLink{sga3:xvi:theorem:1:3}{1.3}. We are reduced to the
case where \(S\) is the spectrum of a discrete valuation ring and
where \(G^0\) possesses a Cartan subgroup \(C\).

\smallskip
\noindent\textit{Proof of \SGARefLink{sga3:xvi:corollary:1:4}{1.4} a).}
We must show that the restriction of \(u\) to \(C\) is an immersion
(\SGARefLink{sga3:xvi:theorem:1:3}{1.3}~a)), which reduces us to the case where \(G=C\) has connected,
nilpotent fibers. Standard reductions (cf.\ the proof of \SGARefLink{sga3:viii:lemma:7:1}{Exposé VIII,
7.1}) reduce us to the case where \(H\) is flat over \(S\), \(u\) is
an isomorphism on the generic fiber, and \(u\) is an isomorphism of
the underlying spaces on the closed fiber. The group \(H\) then has
connected fibers and is therefore separated over \(S\) (Exposé
VI\textsubscript{B}, \SGARefLink{prop:sga3-VIB-5.2}{5.2}). Grant for a moment the following lemma.

\medskip
\noindent\SGARefTarget{sga3:xvi:lemma:1:6}%
\textbf{Lemma 1.6.} --- \textit{Let \(S\) be a prescheme and let
\(G\) be an \(S\)-group prescheme, smooth over \(S\), with connected,
nilpotent fibers and unipotent rank zero. Then:}
\begin{enumerate}
\item[\textit{i)}] \textit{\(G\) is commutative.}
\item[\textit{ii)}] \textit{For every integer \(n>0\),
\({}_{n}G\) is an \(S\)-group prescheme, flat and quasi-finite over
\(S\), and it is finite over \(S\) if and only if the abelian rank,
or the reductive rank, of \(G\) is locally constant on \(S\).}
\item[\textit{iii)}] \textit{For every integer \(q>0\) invertible on
\(S\), the family of subgroups \({}_{q^n}G\) is universally
schematically dense in \(G\) relative to \(S\) (EGA IV, 11.10.8).}
\end{enumerate}

\SGARefLink{sga3:xvi:lemma:1:6}{Lemma 1.6} applies to \(G\), and since \(G\) possesses maximal tori
locally, the reductive rank of the fibers of \(G\) is locally
constant. Thus, by \SGARefLink{sga3:xvi:lemma:1:6}{1.6}~ii), \({}_{n}G\) is finite over \(S\) for
every integer \(n>0\). It follows from \SGARefLink{sga3:viii:proposition:7:9}{Exposé VIII, 7.9} that \(u\) is
an immersion.

\medskip
\noindent\textit{Proof of \SGARefLink{sga3:xvi:lemma:1:6}{1.6}.}
Let us first examine the case where \(S\) is the spectrum of a field.

\medskip
\noindent\SGARefTarget{sga3:xvi:lemma:1:7}%
\textbf{Lemma 1.7.} --- \textit{Let \(k\) be a field of
characteristic \(p\), and let \(G\) be a smooth, connected, nilpotent
algebraic \(k\)-group of unipotent rank zero. Then \(G\) is a
commutative group, an extension of an abelian variety \(A\) by a
torus \(T\). For every integer \(n>0\), \({}_{n}G\) is a finite
group, étale if \((n,p)=1\), defined by a \(k\)-algebra of rank
\(n^{\rho_r+2\rho_{ab}}\). If \((q,p)=1\), the family of subgroups
\({}_{q^n}G\), \(n\in\mathbb N\), is schematically dense in \(G\).}

\smallskip
\noindent\textit{Proof of \SGARefLink{sga3:xvi:lemma:1:7}{1.7}.}
Let \(Z\) be the center of \(G\). The group \(G/Z\) is affine
(\SGARefLink{sga3:xii:6.1}{Exposé XII, 6.1}), smooth, connected, and of unipotent rank zero;
hence it is a torus (BIBLE 4, th.~4). It follows that \(G\) is
commutative (Exposé XII, 6.4). If \(T\) is the unique maximal torus
of \(G\) (Exposé XV, 3.4), Chevalley's theorem (Séminaire Bourbaki
1956/57, no.~145) immediately shows that \(G\) is an extension of an
abelian variety \(A\) by \(T\). For every integer \(n>0\), raising to
the \(n\)-th power is an epimorphism on \(T\), whence the exact
sequence
\[
  0\longrightarrow {}_{n}T\longrightarrow {}_{n}G
  \longrightarrow {}_{n}A\longrightarrow 0.
\]
A classical theorem of Weil (A.~Weil, \emph{Variétés abéliennes et
courbes algébriques}, \S IX, th.~33, cor.~1) says that \({}_{n}A\)
is a finite group defined by a \(k\)-algebra of rank
\(n^{2\rho_{ab}}\). Since \({}_{n}T\) is a finite group of rank
\(n^{\rho_r}\), the asserted structure of \({}_{n}G\) follows.

Now let \(H\) be the smallest closed subscheme of \(G\) containing
\({}_{q^n}G\) for every \(n\). It follows from the preceding
discussion and \SGARefLink{sga3:xv:lemma:4:6}{Exposé XV, 4.6} that \(H\) is a smooth, connected
subgroup, hence of the same type as \(G\). Raising to the \(q\)-th
power on \(H\) is an epimorphism (since \({}_{q}H\) is finite), so
there is an exact sequence
\[
  0\longrightarrow {}_{q}H\longrightarrow {}_{q}G
  \longrightarrow {}_{q}(G/H)\longrightarrow 0.
\]
It follows that \({}_{q}(G/H)=0\), and therefore \(G/H=0\). Thus the
subgroups \({}_{q^n}G\) are schematically dense in \(G\).

\smallskip
\noindent\textit{Continuation of the proof of \SGARefLink{sga3:xvi:lemma:1:6}{1.6}~i).}
To show that \(G\) is commutative, the usual procedure reduces us
first to the case where \(S\) is noetherian, and then to the case
where \(S\) is the spectrum of a local ring with closed point \(s\).
The center of \(G\) is represented by a closed group subscheme \(Z\)
of \(G\) (\SGARefLink{sga3:xi:corollary:6:11}{Exposé XI, 6.11}). To prove that \(Z=G\), it suffices to
prove this after reduction modulo every power of the maximal ideal of
the ring of \(S\): \(Z\) will then be an open subgroup of \(G\), and
hence equal to \(G\), since \(G\) has connected fibers. This reduces
us to the case where \(S\) is local Artinian.

Let \(q\) be an integer invertible on \(S\). For every integer \(n\),
\({}_{q^n}G_s\) is an étale, central subgroup of multiplicative type
of \(G_s\) (\SGARefLink{sga3:xvi:lemma:1:7}{1.7}). Since \(G\) is smooth,
\({}_{q^n}(G_s)\) lifts to an étale, central \(S\)-group subscheme
\(M_n\) of \(G\) (\SGARefLink{sga3:xv:proposition:1:2}{Exposé XV, 1.2}). The family of flat subgroups
\(M_n\), \(n\in\mathbb N\), is then schematically dense in \(G\)
(1.7 and EGA IV, 11.10.9). Since \(Z\) is closed in \(G\) and
contains every \(M_n\), we have \(Z=G\).

\smallskip
\noindent\textit{Proof of \SGARefLink{sga3:xvi:lemma:1:6}{1.6}~ii).}
To see that \({}_{n}G\) is flat and quasi-finite over \(S\), it
suffices to show that the \(n\)-th power map on \(G\) is flat and
quasi-finite. Since \(G\) is flat over \(S\), we are reduced to the
case where \(S\) is the spectrum of a field (EGA IV, 11.3.10), and
we observed in the proof of 1.7 that the \(n\)-th power map is an
epimorphism. Moreover, \({}_{n}G\) is separated over \(S\), since
\(G\) is separated over \(S\) (Exposé VI\textsubscript{B}, \SGARefLink{prop:sga3-VIB-5.2}{5.2}).
For \({}_{n}G\) to be finite over \(S\), it is necessary and
sufficient that the fibers of \({}_{n}G\) be spectra of finite
algebras of locally constant rank. This is seen immediately by
reducing to the case where \(S\) is the spectrum of a henselian local
ring. By \SGARefLink{sga3:xvi:lemma:1:7}{1.7}, this condition is equivalent to saying that the abelian
rank, or the reductive rank, of the fibers of \(G\) is locally
constant.

\smallskip
\noindent\textit{Proof of \SGARefLink{sga3:xvi:lemma:1:6}{1.6}~iii).}
To see that the family of subgroups
\({}_{q^n}G\), \(n\in\mathbb N\), is universally schematically dense
in \(G\), we again reduce to the noetherian case. In view of 1.7 and
1.6~ii), it suffices to apply EGA IV, 11.10.9.

\medskip
\noindent\textit{Proof of \SGARefLink{sga3:xvi:corollary:1:4}{1.4} b).}
We have reduced to the case where \(S\) is the spectrum of a discrete
valuation ring and where \(G^0\) possesses a Cartan subgroup \(C\).
Let
\[
  N=\Norm_G(C)=\Norm_G(T),
\]
where \(T\) is the unique maximal torus of \(C\) (\SGARefLink{sga3:xii:7.1}{Exposé XII,
7.1}~a) and~b)). Since \(H\) is separated over \(S\), it follows from
\SGARefLink{sga3:xvi:lemma:1:6}{1.6}~ii) and \SGARefLink{sga3:viii:proposition:7:12}{Exposé VIII, 7.12} that the restriction of \(u\) to \(C\)
is a closed immersion. On the other hand, to prove that \(u\) is a
closed immersion, it suffices to prove this for \(u|_N\) (\SGARefLink{sga3:xvi:theorem:1:3}{1.3}~b)).
Since, by hypothesis, \(W=N/C\) is represented by a finite group
scheme, the assertion follows from the next lemma, applied to the
exact sequence
\[
  0\longrightarrow C\longrightarrow N\longrightarrow W
  \longrightarrow 1.
\]
(Recall that a proper immersion is a closed immersion.)

\medskip
\noindent\SGARefTarget{sga3:xvi:lemma:1:8}%
\textbf{Lemma 1.8.} --- \textit{Let \(S\) be a prescheme and let
\(G\) be an \(S\)-group prescheme of finite presentation over \(S\),
an extension of an \(S\)-group prescheme \(G''\), proper and of finite
presentation over \(S\), by an \(S\)-group prescheme \(G'\), flat and
of finite presentation over \(S\):}
\[
  1\longrightarrow G'\longrightarrow G\longrightarrow G''
  \longrightarrow 1.
\]
\textit{Let \(H\) be an \(S\)-group prescheme of finite presentation
over \(S\) (or locally of finite presentation if \(S\) is locally
noetherian), and let \(u:G\to H\) be an \(S\)-group morphism. If the
restriction of \(u\) to \(G'\) is proper, then \(u\) is proper.}

\smallskip
\noindent\textit{Proof of \SGARefLink{sga3:xvi:lemma:1:8}{1.8}.}
As usual, we reduce to the case where \(S\) is noetherian (Exposé
VI\textsubscript{B}, \S10, and Exposé XV, 6.2), and we may then apply
the valuative criterion of properness (EGA II, 7.3.8). We therefore
suppose that \(S\) is the spectrum of a discrete valuation ring
\(A\), with closed point \(s\) and generic point \(t\). Let
\(\bar x\in G(t)\) and \(h\in H(S)\) satisfy
\[
  u_t(\bar x)=h(t).
\]
We must show that \(\bar x\) comes from a unique element
\(x\in G(S)\). It is enough to prove the existence and uniqueness of
\(x\) after a faithfully flat extension of the discrete valuation
ring \(A\). Let \(\bar y\) be the image of \(\bar x\) in \(G''(t)\).
Since \(G''\) is proper over \(S\), \(\bar y\) comes from a unique
element \(y\in G''(S)\). Let \(X\) be the inverse image of \(y\) in
\(G\). The prescheme \(X\) is faithfully flat over \(S\), since
\(G'\) is faithfully flat over \(S\), as is the morphism
\(G\to G''\) (Exposé VI\textsubscript{B}, \S9). After a faithfully
flat change of discrete valuation ring, we may suppose that \(X\)
has a section \(\ell\) over \(S\) (EGA IV, 14.5.8). Replacing
\(\bar x\) by \(\bar x\ell_t^{-1}\), we may suppose that
\(\bar x\in G'(t)\). The existence and uniqueness of \(x\) now
follow from the properness of \(u|_{G'}\).

\medskip
\noindent\textit{Proof of \SGARefLink{sga3:xvi:corollary:1:5}{1.5}.}

\smallskip
\noindent\textit{a)}
We shall see in Exposé XIX that, if \(G\) is reductive, then \(G\)
possesses maximal tori locally for the faithfully flat topology, has
a finite Weyl group, and has unipotent rank zero. Assertion~a)
therefore follows from \SGARefLink{sga3:xvi:corollary:1:4}{1.4}.

\smallskip
\noindent\textit{b)}
If \(G\) has affine fibers and unipotent rank zero, then \(G\)
possesses maximal tori locally for the fpqc topology (\SGARefLink{sga3:xv:corollary:8-18}{Exposé XV,
8.18}). It remains only to apply \SGARefLink{sga3:xvi:corollary:1:4}{1.4}, observing that, since \(G\) has
connected, smooth, affine, solvable fibers, its Weyl group is the
identity group (BIBLE 6, th.~1, cor.~3).

\medskip
\noindent\textit{Proof of \SGARefLink{sga3:xvi:theorem:1:3}{1.3} a).}
Since \(u\) is already a monomorphism, to prove that \(u\) is an
immersion it suffices to prove that \(u\) is proper at the points of
\(u(G)\) (EGA IV, 15.7.1); for this we may use the valuative
criterion of local properness (EGA IV, 15.7.5). We are therefore
reduced to the case where \(S\) is the spectrum of a complete
discrete valuation ring with algebraically closed residue field,
closed point \(s\), and generic point \(t\). Since \(G\) is flat over
\(S\), standard reductions (cf.\ the proof of \SGARefLink{sga3:viii:lemma:7:1}{Exposé VIII, 7.1})
reduce us to the case where \(H\) is flat over \(S\), \(u_t\) is an
isomorphism, and \(u_s\) is an isomorphism of underlying spaces. To
say that \(u\) is an immersion is then equivalent to saying that
\(u\) is an isomorphism.

By hypothesis, \(G\) possesses Cartan subgroups locally for the fpqc
topology, so we may speak of the open subset \(U\) of regular points
of \(G\) (\SGARefLink{sga3:xv:theorem:7-3}{Exposé XV, 7.3}~i) and~ii)).

\medskip
\noindent\SGARefTarget{sga3:xvi:lemma:1:9}%
\textbf{Lemma 1.9.} --- \textit{Under the preceding hypotheses, for
\(u\) to be an immersion it suffices that \(u|_U\) be an immersion.}

Indeed, to say that \(u|_U\) is an immersion means that there exist
an open subscheme \(V\) of \(H\) and a closed subscheme \(F\) of
\(V\) such that \(u|_U\) factors through \(F\) and induces an
isomorphism from \(U\) onto \(F\). Since \(H_t\), and therefore
\(V_t\), is irreducible, one necessarily has \(V_t=F_t\). Then \(F\)
contains the schematic closure of \(V_t\) in \(V\), which equals
\(V\) because \(H\) is flat over \(S\). Thus \(V=F\). It follows
that \(u|_U\) is an open immersion, and hence that \(u\) is an open
immersion (Exposé VI\textsubscript{B}, \SGARefLink{cor:sga3-VIB-2.6}{2.6}).

It remains to prove that \(u|_U\) is an immersion. For this we apply
the valuative criterion of local properness. After replacing \(S\)
by the spectrum of a discrete valuation ring faithfully flat over
\(S\), we must show the following: if \(h\) is a section of \(H\)
over \(S\) whose image \(h(S)\) is contained in \(u(U)\), then \(h\)
is the image of a section \(g\) of \(U\) over \(S\). It suffices to
show that \(h\) is contained in the image of a Cartan subgroup \(C\)
of \(G\). Indeed, by hypothesis \(u|_C\) is an immersion, so \(h\)
is the image of a section \(g\) of \(C\), necessarily a section of
\(U\).

Let \(a\in U(s)\) and \(b\in U(t)\) satisfy, respectively,
\[
  u_s(a)=h(s),
  \qquad
  u_t(b)=h(t).
\]
Since \(u_t\) is an isomorphism, \(h(t)\) is a regular point of
\(H_t\), and is therefore contained in a unique Cartan subgroup
\(D_t\) of \(H_t\) (\SGARefLink{rem:sga3-VIB-3.2}{Exposé XIII, 3.2}). Let \(D\) be the schematic
closure of \(D_t\) in \(H\).

\medskip
\noindent\SGARefTarget{sga3:xvi:lemma:1:10}%
\textbf{Lemma 1.10.} ---
\begin{enumerate}
\item[\textit{i)}] \textit{\(D_s\) is a nilpotent algebraic group
(Exposé VI\textsubscript{B}, \S8).}
\item[\textit{ii)}] \textit{\(\dim D_s=\nu\), where \(\nu\) is the
nilpotent rank of \(G\), and also the nilpotent rank of
\((H_s)_{\mathrm{red}}\).}
\item[\textit{iii)}] \textit{\(h(s)\in D_s(\kappa(s))\).}
\end{enumerate}

\noindent\textit{i)}
Since \(H\) is separated (Exposé VI\textsubscript{B}, \SGARefLink{prop:sga3-VIB-5.2}{5.2}), so is
\(D\). Moreover, \(D\) is flat over \(S\) and its generic fiber is
nilpotent. It follows from Exposé VI\textsubscript{B}, \SGARefLink{prop:sga3-VIB-8.4}{8.4} that
\(D_s\) is nilpotent.

\smallskip
\noindent\textit{ii)}
Since \((H_s)_{\mathrm{red}}\simeq G_s\), these two groups have the
same nilpotent rank. Since \(G\) possesses Cartan subgroups locally
for the fpqc topology, \(G_s\) and \(G_t\) have the same nilpotent
rank \(\nu\). Finally, since \(D\) is flat over \(S\),
\[
  \dim D_s=\dim D_t=\nu
\]
(Exposé VI\textsubscript{B}, \S4).

\smallskip
\noindent\textit{iii)}
Since \(h_t\) factors through \(D_t\), the section \(h\) plainly
factors through \(D\), which proves iii).

The next lemma now shows that \((D_s)_{\mathrm{red}}\) is a Cartan
subgroup of
\[
  (H_s)_{\mathrm{red}}\simeq G_s.
\]

\medskip
\noindent\SGARefTarget{sga3:xvi:lemma:1:11}%
\textbf{Lemma 1.11.} --- \textit{Let \(G\) be a smooth, connected
algebraic group over an algebraically closed field \(k\), and let
\(D\) be a smooth, nilpotent algebraic subgroup containing a regular
element \(a\in G(k)\). Then \(D^0\) is contained in a Cartan subgroup
of \(G\). If, in addition, \(\dim D\) equals the nilpotent rank of
\(G\), then \(D\) is a Cartan subgroup of \(G\), and in particular is
connected.}

Let \(Z\) be the center of \(G\), let \(G'=G/Z\), which is affine
(\SGARefLink{sga3:xii:6.1}{Exposé XII, 6.1}), and let \(a'\) and \(D'\) be the images of \(a\)
and \(D\) in \(G'\). The correspondence between Cartan subgroups of
\(G\) and Cartan subgroups of \(G'\) (\SGARefLink{sga3:xii:6.6}{Exposé XII, 6.6}~e)) shows at
once that \(a'\) is a regular element of \(G'\) and that it suffices
to prove the lemma for the pair \((D',G')\). Thus we may suppose that
\(G\) is affine.

Let \(s\) be the semisimple component of \(a\). It belongs to
\(D(k)\) (BIBLE 4, th.~3) and is regular (BIBLE 7, th.~2, cor.~1).
By BIBLE 6, th.~2, \(s\) centralizes \(D^0\); hence \(D^0\) is
contained in the connected centralizer of \(s\), which is a Cartan
subgroup of \(G\) (BIBLE 7, th.~2). If now
\(\dim D\) equals the nilpotent rank of \(G\), then \(D^0\) is a
Cartan subgroup of \(G\), equal to \(\Centr_G(T)\), where \(T\) is
the unique maximal torus of \(D^0\) (\SGARefLink{sga3:xii:6.6}{Exposé XII, 6.6}). But \(D\) is
nilpotent, and therefore centralizes \(T\) (BIBLE 6, th.~2), so
\(D=D^0\).

Return now to the group \(G\). The group
\[
  C_t=u_t^{-1}(D_t)
\]
is the unique Cartan subgroup of \(G_t\) containing \(b\). Let \(C\)
be the schematic closure of \(C_t\) in \(G\).

By functoriality of schematic closure, \(u|_C\) factors through
\(D\), and therefore \(u_s(C_s)\subset D_s\). Since \(u_s\) is an
isomorphism from \(G_s\) onto \((H_s)_{\mathrm{red}}\), since
\[
  \dim C_s=\dim C_t=\nu=\dim D_s
\]
by \SGARefLink{sga3:xvi:lemma:1:10}{1.10}, and since \(D_s\) is connected by \SGARefLink{sga3:xvi:lemma:1:11}{1.11}, \(u_s\) induces an
isomorphism from \((C_s)_{\mathrm{red}}\) onto
\((D_s)_{\mathrm{red}}\). Thus \((C_s)_{\mathrm{red}}\) is a Cartan
subgroup of \(G_s\). The next lemma shows that \(C\) itself is a
Cartan subgroup of \(G\).

\medskip
\noindent\SGARefTarget{sga3:xvi:lemma:1:12}%
\textbf{Lemma 1.12.} --- \textit{Let \(S\) be the spectrum of a
discrete valuation ring, let \(G\) be a smooth \(S\)-group prescheme
of finite type, and let \(C\) be a group subscheme of \(G\), flat
over \(S\), such that the generic fiber \(C_t\) is a Cartan subgroup
of \(G_t\) and the geometric reduced closed fiber
\((C_{\bar s})_{\mathrm{red}}\) is a Cartan subgroup of
\(G_{\bar s}\). Then \(C\) is a Cartan subgroup of \(G\).}

We must prove that \(C\) is smooth over \(S\), and it suffices to
establish this after a faithfully flat extension of the base. The
hypotheses on \(C\) imply that the nilpotent rank of the fibers of
\(G\) is constant; hence we may speak of the open subset \(U\) of
regular points of \(G\) (\SGARefLink{sga3:xv:theorem:7-3}{Exposé XV, 7.3}). Since
\((C_{\bar s})_{\mathrm{red}}\) is a Cartan subgroup of
\(G_{\bar s}\), \(C_{\bar s}\) contains a regular point
\(a_{\bar s}\) of \(G_{\bar s}\). Since \(C\) is flat over \(S\),
after a faithfully flat extension of the base if necessary
(EGA IV, 14.5.8), we may suppose that \(a_{\bar s}\) is an element
of \(G(s)\) which lifts to a section \(a\in C(S)\). Since \(U\) is
open and contains \(a(s)\), it contains \(a\). Let \(C'\) be the
unique Cartan subgroup of \(G\) containing \(a\) (\SGARefLink{sga3:xv:theorem:7-3}{Exposé XV, 7.3}).
Necessarily \(C_t=C'_t\), and consequently \(C'\) and \(C\) coincide
with the identity component of the schematic closure of \(C_t\) in
\(G\) (observe that both \(C'\) and \(C\) are flat over \(S\)).

\smallskip
\noindent\textit{End of the proof of \SGARefLink{sga3:xvi:theorem:1:3}{1.3} a).}
By hypothesis, the restriction of \(u\) to the Cartan subgroup \(C\)
of \(G\) is an immersion. It is clear that \(h(S)\) is contained in
\(u(C)=D\), and hence \(h\) is the image of a section \(g\in C(S)\).
\hfill\(\square\)

\medskip
\noindent\textit{Proof of \SGARefLink{sga3:xvi:theorem:1:3}{1.3} b).}
We shall again use the valuative criterion of local properness
(EGA IV, 15.5) in the case of an immersion (respectively, the
valuative criterion of properness (EGA II, 7.3.6) in the case of a
closed immersion). This reduces us to the case where \(S\) is the
spectrum of a complete discrete valuation ring with algebraically
closed residue field. Let \(s\) be the closed point and \(t\) the
generic point of \(S\). Replacing \(H\) by the schematic closure in
\(H\) of \(u_t(G_t)\), we may suppose that \(H\) is flat over \(S\)
and that \(u_t\) is an isomorphism.

Let \(G^0\) be the connected component of \(G\). By \SGARefLink{sga3:xvi:remark:1:2}{1.2} a), the
restriction \(u|_{G^0}\) is an immersion; denote by \(H^0\) the group
subprescheme of \(H\) that is the image of \(G^0\). Thus
\[
  H_s^0=(H_s)_{\mathrm{red}}^0.
\]

To verify the valuative criterion, we must show that every section
\(h\) of \(H\) over \(S\) such that \(h(S)\) is contained in
\(u(G)\), in the case of an immersion (respectively, every section
\(h\) of \(H\) over \(S\), in the case of a closed immersion), is
the image of a section \(g\) of \(G\) over \(S\).

Let \(C\) be a Cartan subgroup of \(G\), and let
\(D=u(C)\) be the Cartan subgroup of \(H^0\) that is the image of
\(C\). The \(S\)-group
\[
  D'=\operatorname{int}(h)D
\]
has connected fibers, so its underlying space is contained in that
of \(H^0\); on the other hand, since it is smooth over \(S\), it is
reduced, and hence \(D'\) is a smooth group subprescheme of \(H^0\).
The fibers of \(D'\) are Cartan subgroups of the fibers of \(H^0\);
therefore \(D'\) is a Cartan subgroup of \(H^0\), and
\[
  C'=u^{-1}(D')
\]
is a Cartan subgroup of \(G\). But
\(\operatorname{Transpstr}_G(C,C')\) is smooth and surjective over
\(S\) (\SGARefLink{sga3:xii:7.1}{Exposé XII, 7.1} b)). Since \(S\) is henselian with
algebraically closed residue field, there exists \(g\in G(S)\) such
that \(\operatorname{int}(g)C=C'\). Replacing \(h\) by
\(u(g)^{-1}h\), we may suppose that \(h(t)\) normalizes \(D_t\)
(and that \(h(s)\) is the image of an element of the normalizer of
\(C_s\) in \(G_s\), in the case of an immersion). Let
\[
  N=\Norm_G(C).
\]
By hypothesis, \(u|_N\) is an immersion (respectively, a closed
immersion); hence \(h\) is indeed the image of a section of \(N\)
over \(S\), which completes the proof.
\hfill\(\square\)

% END inputs/XVI/tex/components/01_expose_XVI_section1_immersion_criterion.tex

% BEGIN inputs/XVI/tex/components/02_expose_XVI_section2_quotient_representability.tex
% SGA 3, Exposé XVI, Loop-1 English reconstruction.
% Current scope: Section 2 from its opening through the definition of the
% functor F in the proof of Corollary 2.4.
% Authority: Polo--Gille Expo16.pdf, local pages 10--12
% (combined-reader pages 948--950; printed pages 503--507).
% The two diagrams are source-PDF crops used as Loop-1 placeholders; native
% mathematical reconstruction remains a Loop-2 obligation.

\SGARefTarget{sga3:xvi:section:2}%
\section*{2. A Representability Theorem for Quotients}
\addcontentsline{toc}{section}{2. A Representability Theorem for Quotients}

Let us ``recall'' the following result.

\medskip
\noindent\SGARefTarget{sga3:xvi:theorem:2:1}%
\textbf{Theorem 2.1.} --- \textit{Let \(S\) be a prescheme, let
\(X\) and \(Y\) be two \(S\)-preschemes, and let
\(f:X\to Y\) be an \(S\)-morphism. Suppose that one is in one of the
following two cases:}

\begin{enumerate}
\renewcommand{\labelenumi}{\textit{\alph{enumi})}}
\item \textit{The morphism \(f\) is locally of finite presentation.}

\item \textit{The prescheme \(S\) is locally noetherian and \(X\) is
locally of finite type over \(S\).}
\end{enumerate}

\noindent\textit{Then the following conditions are equivalent:}

\begin{enumerate}
\renewcommand{\labelenumi}{\textit{\roman{enumi})}}
\item \textit{There exist an \(S\)-prescheme \(X'\) and a
factorization of \(f\)}
\[
  f:X\xrightarrow{\,f'\,}X'\xrightarrow{\,f''\,}Y,
\]
\textit{where \(f'\) is a faithfully flat \(S\)-morphism, locally
of finite presentation, and \(f''\) is a monomorphism.}

\item \textit{The first projection}
\[
  p_1:X\times_Y X\longrightarrow X
\]
\textit{is a flat morphism.}
\end{enumerate}

\noindent\textit{Moreover, if the preceding conditions hold, then
\((X',f')\) is a quotient of \(X\), for the fpqc topology, by the
equivalence relation defined by \(f\). Consequently the
factorization \(f=f''\circ f'\) in i) is unique up to isomorphism.}

The proof of this delicate theorem will be given in EGA V. One may
also consult the lecture by J.-P. Murre, \textit{Séminaire
Bourbaki}, May 1965, no.~294, theorem~2, corollary~2, which treats
the case where \(Y\) is locally noetherian and \(X\) is of finite
type over \(Y\). We shall see that one can reduce to that case.

Let us first make a few remarks.

\medskip
\noindent\textit{a)} The implication i) \(\Rightarrow\) ii) is
trivial. Indeed, the first projection
\[
  p'_1:X\times_{X'}X\longrightarrow X
\]
factors through \(X\times_Y X\):
\[
  p'_1:X\times_{X'}X\xrightarrow{\,u\,}
  X\times_Y X\xrightarrow{\,p_1\,}X.
\]
The morphism \(u\) is an isomorphism because \(f''\) is a
monomorphism, and \(p'_1\) is flat because \(f'\) is flat; therefore
\(p_1\) is flat.

\medskip
\noindent\textit{b)} The assertions of \SGARefLink{sga3:xvi:theorem:2:1}{2.1} are local on \(Y\)
(hence local on \(S\)); they are also local on \(X\), as follows
readily from the fact that a flat morphism locally of finite
presentation is open (EGA IV, 11.3.1).

\medskip
\noindent\textit{c)} Under the hypotheses of \SGARefLink{sga3:xvi:theorem:2:1}{2.1} a), the foregoing
reduces us to the case where \(X\) and \(Y\) are affine and \(f\) is
of finite presentation. Replacing \(S\) by \(Y\), we may suppose
that \(X\) and \(Y\) are of finite presentation over \(S\). We then
reduce to the noetherian case for \(S\) by EGA IV, 11.2.6.

\medskip
\noindent\textit{d)} Under the hypotheses of \SGARefLink{sga3:xvi:theorem:2:1}{2.1} b), we may suppose
that \(S\), \(X\), and \(Y\) are affine, that \(S\) is noetherian,
and that \(X\) is of finite type over \(S\). Regard \(Y\) as a
filtered projective limit of affine schemes \(Y_i\) of finite type
over \(S\). The schemes \(X\times_{Y_i}X\) form a decreasing
filtered family of closed subschemes of \(X\times_S X\), whose
projective limit is \(X\times_Y X\). Since \(X\times_S X\) is
noetherian, one has
\[
  X\times_{Y_i}X=X\times_Y X
\]
for all sufficiently large \(i\). Thus
\[
  f_i:X\xrightarrow{\,f\,}Y\longrightarrow Y_i
\]
satisfies the hypotheses of \SGARefLink{sga3:xvi:theorem:2:1}{2.1} ii) if \(f\) does. Since the
equivalence relation on \(X\) defined by \(f\) coincides with the
one defined by \(f_i\), it is clear that it suffices to prove
ii) \(\Rightarrow\) i) for \(f_i\). This reduces us to the case
where \(Y\) is of finite type over \(S\).

\subsection*{Application to Group Preschemes}
\addcontentsline{toc}{subsection}{Application to Group Preschemes}

\noindent\SGARefTarget{sga3:xvi:theorem:2:2}%
\textbf{Theorem 2.2.} --- \textit{Let \(S\) be a prescheme and let
\(G\) be an \(S\)-group prescheme locally of finite presentation
over \(S\), acting on an \(S\)-prescheme \(X\). Let
\(\xi\in X(S)\) be a section of \(X\) over \(S\), such that the
stabilizer \(H\) of \(\xi\) in \(G\) is an \(S\)-group
subprescheme of \(G\), flat over \(S\). If \(X\) is locally of
finite type over \(S\), or if \(S\) is locally noetherian, then the
homogeneous quotient space \(G/H\) is representable by an
\(S\)-prescheme locally of finite presentation over \(S\), and the
\(S\)-morphism}
\[
  f:G\longrightarrow X,\qquad g\longmapsto g\cdot\xi
\]
\textit{factors as}
\begin{center}
  % BEGIN inputs/XVI/native_diagrams/exp16/001_theorem_2_2_factorization_triangle.tex
\begin{tikzcd}[column sep=5.6em,row sep=2.6em]
  G \arrow[d,"p"'] \arrow[dr,"f"]
    & {} \\
  G/H \arrow[r,"i"]
    & X.
\end{tikzcd}

% END inputs/XVI/native_diagrams/exp16/001_theorem_2_2_factorization_triangle.tex

\end{center}
\textit{where \(p\) is the canonical projection, a faithfully flat
morphism locally of finite presentation, and \(i\) is a
monomorphism. (For definiteness, we have supposed that \(G\) acts
on the left on \(X\).)}

\medskip
\noindent\textit{Proof.}
The morphism \(f\) makes \(G\) into an \(X\)-prescheme. By the
definition of the stabilizer of \(\xi\), the morphism
\[
  G\times_S H\longrightarrow G\times_X G,\qquad
  (g,h)\longmapsto(g,gh)
\]
is an isomorphism. Since \(H\) is flat over \(S\),
\(G\times_S H\) is flat over \(G\); hence the first projection
\[
  p_1:G\times_X G\longrightarrow G
\]
is flat. Moreover, if \(X\) is locally of finite type over \(S\),
then \(f\) is locally of finite presentation (EGA IV, 1.4.3 v));
otherwise \(S\) is locally noetherian by hypothesis. It therefore
suffices to apply \SGARefLink{sga3:xvi:theorem:2:1}{2.1} to \(f\).

It remains only to show that \(G/H\) is locally of finite
presentation over \(S\), and this follows immediately from
\SGARefLink{sga3:v:proposition:9:1}{Exposé V, 9.1}.
\hfill\(\square\)

\medskip
\noindent\SGARefTarget{sga3:xvi:corollary:2:3}%
\textbf{Corollary 2.3.} --- \textit{Let \(S\) be a prescheme, let
\(G\) and \(H\) be two \(S\)-group preschemes, and let
\(u:G\to H\) be an \(S\)-homomorphism. Suppose that \(G\) is locally
of finite presentation over \(S\), and that either \(H\) is locally
of finite type over \(S\), or \(S\) is locally noetherian. If
\(K=\Ker(u)\) is flat over \(S\), then the quotient group \(G/K\)
is representable by an \(S\)-group prescheme locally of finite
presentation over \(S\), and \(u\) factors as}
\begin{center}
  % BEGIN inputs/XVI/native_diagrams/exp16/002_corollary_2_3_factorization_triangle.tex
\begin{tikzcd}[column sep=5.8em,row sep=4.0em]
  G \arrow[rr,"u"] \arrow[dr,"p"']
    & {}
    & H \\
  {}
    & G/K \arrow[ur,"i"']
    & {}
\end{tikzcd}

% END inputs/XVI/native_diagrams/exp16/002_corollary_2_3_factorization_triangle.tex

\end{center}
\textit{where \(p\) is the canonical projection and \(i\) is a
monomorphism.}

\medskip
\noindent\textit{Proof.}
Apply \SGARefLink{sga3:xvi:theorem:2:2}{2.2} with \(X=H\) and with \(\xi\) the unit section of \(H\).
\hfill\(\square\)

\medskip
\noindent\SGARefTarget{sga3:xvi:corollary:2:4}%
\textbf{Corollary 2.4.} --- \textit{Let \(S\) be a prescheme, let
\(G\) be an \(S\)-group prescheme of finite presentation over
\(S\), let \(H\) be an \(S\)-group prescheme smooth over \(S\) with
connected fibers (and hence of finite presentation over \(S\), by
Exposé VI\textsubscript{B}, \SGARefLink{cor:sga3-VIB-5.5}{5.5}), and let \(i:H\to G\) be a
monomorphism of \(S\)-groups, so that \(H\) is a subgroup of \(G\).
Suppose that}
\[
  N=\Norm_G(H),
\]
\textit{which is representable by a closed group subprescheme of
\(G\), of finite presentation over \(S\) (\SGARefLink{sga3:xi:corollary:6:11}{Exposé XI, 6.11}), is flat
over \(S\). Then \(G/N\) is representable by an \(S\)-prescheme of
finite presentation over \(S\) and quasi-projective over \(S\).}

All the assertions to be proved, except the last one, are local on
\(S\). To establish them, we may therefore suppose that \(S\) is
quasi-compact and that the relative dimension of \(H\) over \(S\)
is constant, equal to \(r\). Proceed as in Exposé XV, \S5. For every
integer \(n>0\), let \(G^{(n)}\) (respectively \(H^{(n)}\)) be the
\(n\)-th normal invariant of the unit section of \(G\)
(respectively of \(H\)) (EGA IV, 16). The sheaf of
\(\mathcal O_S\)-modules \(H^{(n)}\) is a quotient of \(G^{(n)}\).
Since \(H\) is smooth over \(S\) of dimension \(r\), \(H^{(n)}\)
canonically defines an element \(\xi_n\) of
\[
  X_n=\operatorname{Grass}_{\varphi(n,r)}G^{(n)}
\]
(EGA I, second edition, \S9), for a suitable integer
\(\varphi(n,r)\). On the other hand, \(G\) acts naturally on
\(G^{(n)}\), and hence also on \(X_n\), through the representation
\[
  G\longrightarrow\operatorname{Aut}_{S\text{-gr}}(G),
  \qquad g\longmapsto\operatorname{int}(g).
\]
Since \(S\) is quasi-compact, there exists an integer \(m\) such
that, for \(n\geq m\),
\[
  N=\Norm_G(H)=\Norm_G H^{(n)}
  \qquad\text{(Exposé XI, 6.11 b)).}
\]
In other words, \(N\) is the stabilizer of \(\xi_n\) for sufficiently
large \(n\). The representability of \(G/N\) therefore follows from
\SGARefLink{sga3:xvi:theorem:2:2}{2.2}. The fact that \(G/N\) is of finite presentation over \(S\) is
a consequence of \SGARefLink{sga3:v:proposition:9:1}{Exposé V, 9.1}. It remains to prove that \(G/N\) is
quasi-projective over \(S\). For this, we shall exhibit a canonical
invertible sheaf on \(G/N\) that is ample relative to \(S\).
Consider the functor
\[
  \mathscr F:(\mathbf{Sch}/S)^\circ\longrightarrow\mathbf{Ens}
\]
such that, for every \(S\)-prescheme \(T\),
\[
  \mathscr F(T)=
  \left\{
  \begin{array}{c}
  \text{\(T\)-subgroups \(H\) of \(G_T\), representable and smooth
  over \(T\),}\\
  \text{with connected fibers}
  \end{array}
  \right\}.
\]

Such a group \(H\) is of finite presentation over \(S\)
(Exposé VI\textsubscript{B}, 5.5), and
\(H\to G_T\) is a quasi-affine morphism (EGA IV, 8.11.2).
By effective descent for quasi-affine morphisms (SGA 1, VIII, 7.9),
it follows that \(\mathscr F\) is a sheaf for the fpqc topology.
Since \(G/N\) is the sheaf associated with a subfunctor of
\(\mathscr F\), there is a canonical monomorphism
\[
  G/N\longrightarrow\mathscr F.
\]
Consequently there exists a subgroup \(H'\) of
\(G\times_S(G/N)\), representable and smooth over \(G/N\), with
connected fibers, that is ``universal'' for the functor \(G/N\).
I claim that the invertible sheaf
\[
  \mathscr L=\bigl(\det(\operatorname{Lie}H')\bigr)^{-1}
\]
is ample relative to \(S\). In this form the assertion is local on
\(S\), and its proof is analogous to the one given in \SGARefLink{sga3:xv:remark:5:8}{Exposé XV,
5.8}.

The following corollary was announced in \SGARefLink{sga3:xiv:4.8.1}{Exposé XIV, 4.8 bis}.

\medskip
\noindent\SGARefTarget{sga3:xvi:corollary:2:5}%
\textbf{Corollary 2.5.} --- \textit{Let \(S\) be a prescheme, let
\(G\) be an \(S\)-group prescheme smooth and of finite presentation
over \(S\), with connected fibers, and let \(P\) be a parabolic
subgroup of \(G\) (\SGARefLink{sga3:xiv:paragraph:4-8-bis}{Exposé XIV, 4.8 bis}). Then \(G/P\) is
representable by an \(S\)-prescheme smooth and projective over
\(S\).}

It remains only to show (\textit{loc.\ cit.}) that \(G/P\) is
representable by an \(S\)-prescheme that is quasi-projective and of
finite presentation. This follows from \SGARefLink{sga3:xvi:corollary:2:4}{2.4}, since
\[
  P=\Norm_G(P)
\]
(\textit{loc.\ cit.}) and is therefore flat over \(S\).
\hfill\(\square\)

% END inputs/XVI/tex/components/02_expose_XVI_section2_quotient_representability.tex

% BEGIN inputs/XVI/tex/components/03_expose_XVI_section3_flat_center.tex
% SGA 3, Exposé XVI, Loop-1 English reconstruction.
% Current scope: all of Section 3, including the terminal counterexample.
% Authority: Polo--Gille Expo16.pdf, local pages 13--17
% (combined-reader pages 951--955; printed pages 509--518).
% The Proposition 3.1 diagram is a source-PDF crop used as a Loop-1
% placeholder; native mathematical reconstruction remains a Loop-2 task.

\SGARefTarget{sga3:xvi:section:3}%
\section*{3. Groups with Flat Center}
\addcontentsline{toc}{section}{3. Groups with Flat Center}

\noindent\SGARefTarget{sga3:xvi:proposition:3:1}%
\textbf{Proposition 3.1.} --- \textit{Let \(S\) be a prescheme and
let \(G\) be an \(S\)-group prescheme smooth and of finite
presentation over \(S\), with connected fibers. Suppose that the
center \(Z\) of \(G\), which is representable by \SGARefLink{sga3:xi:corollary:6:11}{Exposé XI, 6.11},
is flat over \(S\). For every integer \(n>0\), let \(G^{(n)}\)
denote the locally free \(\mathcal O_S\)-module equal to the
\(n\)-th normal invariant of \(G\) along the unit section, and let
\(\rho_n\) be the natural ``adjoint representation'' of \(G\) in
\(\operatorname{GL}_S(G^{(n)})\). Then:}

\begin{enumerate}
\renewcommand{\labelenumi}{\textit{\alph{enumi})}}
\item \textit{The quotient \(G'=G/Z\) is representable by an
\(S\)-group prescheme that is smooth, of finite presentation over
\(S\), quasi-affine, and has affine connected fibers.}

\item \textit{The representation \(\rho_n\) factors through \(G'\):}
\begin{center}
  % BEGIN inputs/XVI/native_diagrams/exp16/003_proposition_3_1_factorization_triangle.tex
\begin{tikzcd}[column sep=6.0em,row sep=4.0em]
  G \arrow[rr,"\rho_n"] \arrow[dr]
    & {}
    & \operatorname{GL}_S\!\left(G^{(n)}\right) \\
  {}
    & G/Z \arrow[ur,"i_n"']
    & {}
\end{tikzcd}

% END inputs/XVI/native_diagrams/exp16/003_proposition_3_1_factorization_triangle.tex

\end{center}
\textit{If \(S\) is quasi-compact, then \(i_n\) is a monomorphism
for all sufficiently large \(n\).}

\item \textit{The functor \(\mathscr T_{G'}\) of subtori of \(G'\)
is representable by an \(S\)-prescheme smooth over \(S\), and, if
\(S\) is quasi-compact, this prescheme is a sum of a family of
preschemes affine over \(S\).}
\end{enumerate}

\medskip
\noindent\textit{Proof.}
The assertions in a) are local on \(S\), which allows us to suppose
that \(S\) is quasi-compact. By \SGARefLink{sga3:xi:corollary:6:11}{Exposé XI, 6.11} b), one has
\[
  Z=\Ker\rho_n
\]
for all sufficiently large \(n\); this kernel is therefore flat
over \(S\). The representability of \(G/Z\) and the fact that \(i_n\)
is a monomorphism now follow from \SGARefLink{sga3:xvi:corollary:2:3}{2.3}. Since \(G\) is smooth and of
finite presentation over \(S\), the same is true of \(G'\)
(Exposé VI\textsubscript{B}, \S9). The group
\(\operatorname{GL}_S(G^{(n)})\) is affine over \(S\), and a
monomorphism of finite presentation is quasi-affine
(EGA IV, 8.11.2). Hence \(G'\) is quasi-affine over \(S\) and has
affine fibers. Since \(G'\) is smooth over \(S\), the functor
\(\mathscr T_{G'}\) is formally smooth (\SGARefLink{sga3:xi:proposition:2:1:bis}{Exposé XI, 2.1 bis}).
Taking \SGARefLink{sga3:xi:remark:4:6}{Exposé XI, 4.6} and \SGARefLink{sga3:xvi:corollary:4:3}{4.3} into account, the assertions of \SGARefLink{sga3:xvi:proposition:3:1}{3.1} c)
will therefore follow from the next lemma.

\medskip
\noindent\SGARefTarget{sga3:xvi:lemma:3:2}%
\textbf{Lemma 3.2.} --- \textit{Let \(S\) be a prescheme, let \(G\)
and \(H\) be two \(S\)-group preschemes of finite presentation over
\(S\), let \(u:G\to H\) be an \(S\)-group monomorphism, and let
\(\mathscr T_G\) (respectively \(\mathscr T_H\)) be the functor of
subtori of \(G\) (respectively of \(H\)) (cf.\ Exposé XV, \S8).
Then the map \(T\mapsto u(T)\) defines a monomorphism
(\SGARefLink{sga3:xv:lemma:8:3}{Exposé XV, 8.3} c))}
\[
  \widetilde u:\mathscr T_G\longrightarrow\mathscr T_H
\]
\textit{that is representable by a closed immersion of finite
presentation.}

By the usual procedure, we reduce to the following problem. Let
\(T\) be a subtorus of \(H\), let
\[
  T'=G\times_H T
\]
be its inverse image in \(G\), and let \(u_T:T'\to T\) be the
\(S\)-group monomorphism induced by \(u\). One must show that the
\(S\)-functor
\[
  \prod_{T/S}(T'/T)
\]
is representable by a closed subprescheme of \(S\), of finite
presentation. But \(T\), being a torus, is smooth over \(S\) with
connected fibers, and it suffices to apply \SGARefLink{sga3:xi:corollary:6:10}{Exposé XI, 6.10}.
\hfill\(\square\)

\medskip
\noindent\SGARefTarget{sga3:xvi:theorem:3:3}%
\textbf{Theorem 3.3.} --- \textit{Let \(S\) be a prescheme, let
\(G\) be an \(S\)-group prescheme smooth and of finite presentation
over \(S\), with connected fibers, and let \(Z\) be the center of
\(G\). Suppose that \(Z\) is flat over \(S\) and that the unipotent
rank of \(G\) (\SGARefLink{sga3:xv:definition:6:1:ter}{Exposé XV, 6.1 ter}) equals that of \(Z\). Then:}

\begin{enumerate}
\renewcommand{\labelenumi}{\textit{\alph{enumi})}}
\item \textit{The group \(G'=G/Z\) is representable; if \(S\) is
quasi-compact, the canonical morphism}
\[
  i_n:G'\longrightarrow\operatorname{GL}_S(G^{(n)})
\]
\textit{(cf.\ \SGARefLink{sga3:xvi:proposition:3:1}{3.1} b)) is an immersion for all sufficiently large
\(n\).}

\item \textit{The group \(G'\) is quasi-affine over \(S\), with
affine fibers; its center is the unit group, and \(G'\) has
unipotent rank zero.}

\item \textit{Locally for the étale topology, \(G'\) possesses
maximal tori, and these are also Cartan subgroups of \(G'\). The
functor \(\mathscr{TM}_{G'}\) of maximal tori of \(G'\)
(Exposé XV, \S8) is representable by an \(S\)-prescheme smooth and
affine over \(S\).}

\item \textit{Locally for the étale topology, \(G\) possesses
Cartan subgroups, and the functor \(\mathscr C_G\) of Cartan
subgroups of \(G\) is representable by an \(S\)-prescheme smooth
and affine over \(S\).}
\end{enumerate}

\medskip
\noindent\textit{Proof.}
By \SGARefLink{sga3:xvi:proposition:3:1}{3.1}, the group \(G'\) is represented by an \(S\)-prescheme that
is smooth and quasi-affine over \(S\), with affine fibers. In view
of the correspondence between Cartan subgroups of the fibers of
\(G\) and those of the fibers of \(G'\) (\SGARefLink{sga3:xii:6.6}{Exposé XII, 6.6} e)), the
hypothesis on the unipotent rank of \(Z\) implies that \(G'\) has
unipotent rank zero. \SGARefLink{sga3:xv:corollary:8-18}{Exposé XV, 8.18} now shows that, locally for the
étale topology, \(G'\) possesses maximal tori. The fact that \(i_n\)
is an immersion for all sufficiently large \(n\) then follows from
\SGARefLink{sga3:xvi:proposition:3:1}{3.1} b) and \SGARefLink{sga3:xvi:corollary:1:4}{1.4} a). This proves a).

We now prove c). Since \(G'\) has unipotent rank zero, every maximal
torus of \(G'\) is plainly also a Cartan subgroup of \(G'\). The
functor of maximal tori of \(G'\) is represented by a subprescheme
that is both open and closed in the functor of subtori of \(G'\);
moreover it is of finite type over \(S\) (\SGARefLink{sga3:xv:theorem:8-15}{Exposé XV, 8.15}).
It follows from \SGARefLink{sga3:xvi:proposition:3:1}{3.1} c) that this functor is represented by an
\(S\)-prescheme smooth and affine over \(S\).

Since, locally for the étale topology, \(G'\) possesses Cartan
subgroups, the same is true of \(G\); and the functor of Cartan
subgroups of \(G\) is canonically isomorphic to that of \(G'\)
(\SGARefLink{sga3:xv:theorem:7-3}{Exposé XV, 7.3} iv)). Thus c) implies d).

It remains to prove b), and more precisely to prove that the center
\(Z'\) of \(G'\) is the unit group. Since \(Z'\) is representable
(\SGARefLink{sga3:xi:corollary:6:11}{Exposé XI, 6.11}) and of finite presentation over \(S\), it suffices
to show that the fibers of \(Z'\) are reduced to the unit group
(Exposé VI\textsubscript{B}, \SGARefLink{lem:sga3-VIB-2.9}{2.9}). We are thereby reduced to the
case where \(S\) is the spectrum of an algebraically closed field.
The center \(Z'\) is evidently contained in every Cartan subgroup
of \(G'\), hence, by c), in every maximal torus of \(G'\); it is
therefore of multiplicative type. On the other hand, we shall see
in Exposé XVII that, if \(G\) is a connected algebraic group with
center \(Z\), the center \(Z'\) of \(G/Z\) is unipotent. In the
present case, \(Z'\), being at once of multiplicative type and
unipotent, is reduced to the unit group (cf.\ Exposé XVII).
\hfill\(\square\)

\subsection*{Examples of Groups Whose Center Is Flat}
\addcontentsline{toc}{subsection}{Examples of Groups Whose Center Is Flat}

\noindent\SGARefTarget{sga3:xvi:proposition:3:4}%
\textbf{Proposition 3.4.} --- \textit{Let \(S\) be a prescheme, let
\(G\) be an \(S\)-group prescheme of finite presentation over
\(S\), smooth, with connected fibers, and let \(Z\) be the center
of \(G\). Then \(Z\) is flat over \(S\) in either of the following
two cases:}

\begin{enumerate}
\renewcommand{\labelenumi}{\textit{\alph{enumi})}}
\item \textit{The unipotent rank \(\rho_u\) of the fibers of \(G\)
(\SGARefLink{sga3:xv:definition:6:1:ter}{Exposé XV, 6.1 ter}) is zero.}

\item \textit{i) \(S\) is reduced; \quad
ii) the dimension of the fibers of \(Z\) is a locally constant
function on \(S\); \quad
iii) the unipotent rank of \(G\) equals that of \(Z\).}
\end{enumerate}

\medskip
\noindent\textit{Proof of \SGARefLink{sga3:xvi:proposition:3:4}{3.4} a).}
We shall prove at the same time the following proposition.

\medskip
\noindent\SGARefTarget{sga3:xvi:proposition:3:5}%
\textbf{Proposition 3.5.} --- \textit{Under the hypotheses of
\SGARefLink{sga3:xvi:proposition:3:4}{3.4} a), suppose moreover that \(G\) possesses a maximal torus \(T\),
and let}
\[
  C=\Centr_G(T)
\]
\textit{be the Cartan subgroup of \(G\) associated with \(T\)
(\SGARefLink{sga3:xii:7.1}{Exposé XII, 7.1}). Then:}

\begin{enumerate}
\renewcommand{\labelenumi}{\textit{(\roman{enumi})}}
\item \textit{\(Z\cap T\) is a subgroup of multiplicative type of
\(G\).}

\item \textit{\(C\) is commutative and equals \(T\cdot Z\).}

\item \textit{If the quotients \(Z/(Z\cap T)\) and \(C/T\) are
representable, then they are represented by abelian preschemes
(that is, \(S\)-group preschemes smooth over \(S\), whose fibers are
abelian varieties), and the canonical monomorphism}
\[
  Z/(Z\cap T)\longrightarrow C/T
\]
\textit{is an isomorphism.}
\end{enumerate}

Using the general properties of passage to the limit established
in Exposé VI\textsubscript{B}, \S10 and \SGARefLink{sga3:xv:proposition:6-2}{Exposé XV, 6.2}, \SGARefLink{sga3:xvi:proposition:6:3}{6.3}, and
\SGARefLink{sga3:xv:corollary:6:3:bis}{6.3 bis}, we reduce as usual to the case where \(S\) is noetherian
(observe that the assertions of \SGARefLink{sga3:xvi:proposition:3:4}{3.4} and \SGARefLink{sga3:xvi:proposition:3:5}{3.5} are local on \(S\)).
We noted in 3.1 that the hypotheses on \(G\) imply that \(Z\) is
representable. To prove that \(Z\) is flat over \(S\), EGA
0\textsubscript{III}, 10.2.6 reduces us to the case where \(S\) is
local artinian. Since \(G\) is then smooth, it possesses maximal
tori locally for the étale topology (\SGARefLink{sga3:xv:corollary:8:17}{Exposé XV, 8.17}). After a
finite flat extension of \(S\), if necessary (which is permissible
when proving that \(Z\) is flat), we may therefore suppose that
\(G\) possesses a maximal torus \(T\).

Let us show, in the same way, that to establish \SGARefLink{sga3:xvi:proposition:3:5}{3.5} we may reduce
to the case where \(S\) is artinian.

\medskip
\noindent i) Since \(Z\) is closed in \(G\) (\SGARefLink{sga3:xi:corollary:6:11}{Exposé XI, 6.11}),
\(Z\cap T\) is a closed group subscheme of \(T\). It follows from
\SGARefLink{sga3:x:corollary:4-8}{Exposé X, 4.8} b), that \(Z\cap T\) is of multiplicative type if and
only if it is flat over \(S\). As above, it suffices to establish
that \(Z\cap T\) is flat when \(S\) is artinian.

\medskip
\noindent ii) Since \(G\) has unipotent rank zero, every Cartan
subgroup \(C\) of \(G\) satisfies the hypotheses of \SGARefLink{sga3:xvi:lemma:1:6}{Lemma 1.6}, and
is therefore commutative. The equality \(C=T\cdot Z\) will follow
from iii).

\medskip
\noindent iii) If \(C/T\) is representable, then \(C/T\) is smooth
over \(S\) (Exposé VI\textsubscript{B}, \S9), and its fibers are
abelian varieties (\SGARefLink{sga3:xvi:lemma:1:7}{1.7}); thus \(C/T\) is an abelian prescheme. To
show that the canonical monomorphism
\[
  i:Z/(Z\cap T)\longrightarrow C/T
\]
is an isomorphism, it suffices to verify this when \(S\) is local
artinian. Indeed, one then deduces successively that
\(Z/(Z\cap T)\) is flat over \(S\) (EGA
0\textsubscript{III}, 10.2.6), that \(i\) is flat
(EGA IV, 11.3.10), that \(i\) is an open immersion
(EGA IV, 17.9.1), and finally that \(i\) is an isomorphism
(because \(C/T\) has connected fibers).

We henceforth suppose that \(S\) is local artinian, with closed
point \(s\), and that \(G\) possesses a maximal torus \(T\). Put
\[
  C=\Centr_G(T).
\]
The group \(C\) is a Cartan subgroup of \(G\) (\SGARefLink{sga3:xii:7.1}{Exposé XII, 7.1})
and contains \(Z\).

The algebraic group \(M_s=Z_s\cap T_s\) is a central subgroup of
multiplicative type of \(G_s\). Since \(T\) is smooth over \(S\),
\(M_s\) lifts to a group subscheme \(M\) of \(T\), of
multiplicative type (\SGARefLink{sga3:ix:theorem:3:6:bis}{Exposé IX, 3.6 bis}), and contained in the
center of \(G\) (\SGARefLink{sga3:ix:corollary:3:9}{Exposé IX, 3.9}), hence contained in \(Z\cap T\).
Since \(S\) is artinian and \(M\) and \(T\) are flat over \(S\),
the quotient groups
\[
  (Z\cap T)/M,\qquad Z'=Z/M,\qquad A=C/T
\]
are representable (Exposé VI\textsubscript{A}, \S\S4 and 5).
By construction, \(\bigl((Z\cap T)/M\bigr)_s\) is the unit group,
so \(Z\cap T=M\) (Exposé VI\textsubscript{B}, \SGARefLink{lem:sga3-VIB-2.9}{2.9}); in
particular it is flat over \(S\), which proves \SGARefLink{sga3:xvi:proposition:3:5}{3.5} i).

Passing to the quotient gives a canonical monomorphism
\[
  i:Z'\longrightarrow A,
\]
which is therefore a closed immersion (Exposé
VI\textsubscript{B}, \SGARefLink{cor:sga3-VIB-1.4.2}{1.4.2}). We have already observed that \(A\)
is an abelian prescheme. It remains to show that \(i\) is an
isomorphism. This will prove \SGARefLink{sga3:xvi:proposition:3:5}{3.5} iii) and will imply that \(Z'\)
is flat over \(S\), and hence that \(Z\) is flat over \(S\)
(as an extension of \(Z'\) by the flat group \(M\); Exposé
VI\textsubscript{B}, \S9). Since \(Z_s\) has the same abelian rank
as \(G_s\) (\SGARefLink{sga3:xii:6.1}{Exposé XII, 6.1}), \(i_s\) is an epimorphism, hence an
isomorphism. Let \(q>0\) be an integer invertible on \(S\). By the
density theorem of \SGARefLink{sga3:xvi:lemma:1:6}{1.6} iii), to see that \(i(Z')=A\), it suffices
to show that, for every integer \(n\) equal to a power of \(q\),
\(i(Z')\) contains \({}_nA\). Now let \(M_s^0\) be the identity
component of \(M_s\). It follows immediately by duality that
raising to the \(q\)-th power in \(M_s^0\) is an epimorphism.
Consequently, if \(m_0\) is the exponent of \(q\) in the prime
factorization of \(\operatorname{card}(M_s/M_s^0)\), the image of
\({}_{nm_0}Z_s\) in
\[
  A_s\simeq Z_s/M_s
\]
contains \({}_nA_s\). There is therefore a subgroup \(M_s(n)\) of
multiplicative type of \(Z_s\) whose image in \(A_s\) is
\({}_nA_s\). As above, \(M_s(n)\) lifts to a central subgroup
\(M(n)\) of multiplicative type of \(G\). The image of \(M(n)\)
in \(A\) is a subgroup of multiplicative type (\SGARefLink{sga3:ix:theorem:6:8}{Exposé IX, 6.8}),
and is necessarily equal to \({}_nA\), since it is so on the
reduced fiber (\SGARefLink{sga3:ix:theorem:5:1:bis}{Exposé IX, 5.1 bis}). This completes the proof of
\SGARefLink{sga3:xvi:proposition:3:4}{3.4} a) and of \SGARefLink{sga3:xvi:proposition:3:5}{3.5}.

\medskip
\noindent\textit{Proof of \SGARefLink{sga3:xvi:proposition:3:4}{3.4} b).}
The assertion to be proved is local on \(S\), so we may suppose
that \(S\) is affine, with ring \(A\). Regarding \(A\) as the
inductive limit of its finitely generated
\(\mathbb Z\)-subalgebras, we reduce as above to the case where
\(S\) is noetherian and reduced.

To prove that \(Z\) is flat, we may use the valuative criterion
for flatness (EGA IV, 11.8.1), which reduces us to the case where
\(S\) is the spectrum of a complete discrete valuation ring with
algebraically closed residue field, with generic point \(t\) and
closed point \(s\). Let \(Z'\) be the schematic closure of \(Z_t\)
in \(Z\). We must show that \(Z'=Z\), and it even suffices to show
that \(Z_s=Z'_s\) (Exposé VI\textsubscript{B}, \SGARefLink{cor:sga3-VIB-2.6}{2.6}). By Exposé
VI\textsubscript{B}, \S4, one has the inequalities
\[
  \dim Z_t=\dim Z'_t=\dim Z'_s\leq \dim Z_s.
\]
By hypothesis, however, \(\dim Z_t=\dim Z_s\); hence
\(\dim Z'_s=\dim Z_s\), and consequently \(Z'_s\) contains
\((Z_s)^0_{\mathrm{red}}\). It follows that
\[
  G'_s=G_s/Z'_s
\]
is affine (\SGARefLink{sga3:xii:6.1}{Exposé XII, 6.1}). In view of the correspondence between
Cartan subgroups of \(G\) and \(G'\) (\SGARefLink{sga3:xii:6.6}{Exposé XII, 6.6} e)),
hypothesis \SGARefLink{sga3:xvi:proposition:3:4}{3.4} b) iii) implies that the unipotent rank of \(G'_s\)
is zero, and consequently its Cartan subgroups are also its
maximal tori. The image \(Z''_s\) of \(Z_s\) in \(G'_s\) is a
central algebraic subgroup of \(G'_s\), hence is contained in
every Cartan subgroup of \(G'_s\), that is, in every maximal
torus; \(Z''_s\) is therefore a subgroup of multiplicative type
of \(G'_s\). Under these conditions, we shall show in Exposé XVII,
\S7, that there exists a finite subgroup \(M_s\) of multiplicative
type of \(Z_s\) whose image in \(Z_s/Z'_s\) is \(Z''_s\) (we used
this fact in the proof of \SGARefLink{sga3:xvi:proposition:3:4}{3.4} a) when \(Z''_s\) is étale). Since
\(G\) is smooth over \(S\) and \(S\) is the spectrum of a complete
local ring, \(M_s\) lifts to an \(S\)-subgroup \(M\) of
multiplicative type of \(G\) (\SGARefLink{sga3:ix:theorem:3:6:bis}{Exposé IX, 3.6 bis} and \SGARefLink{sga3:xv:proposition:1-6}{Exposé XV,
1.6} b)), which is central (\SGARefLink{sga3:ix:theorem:5:6}{Exposé IX, 5.6} a)). Thus \(M_t\) is
contained in \(Z'_t=Z_t\), and, since \(M\) is flat, \(M\) is
contained in \(Z'\). The group \(M_s\) is therefore contained in
\(Z'_s\), but this implies that \(Z''_s\) is the unit group, that
is, \(Z_s=Z'_s\). This completes the proof of \SGARefLink{sga3:xvi:proposition:3:4}{3.4} b).

\subsection*{Example of a Smooth Group Prescheme with Connected Fibers
Whose Center Is Not Flat}
\addcontentsline{toc}{subsection}{Example of a Smooth Group Prescheme
with Connected Fibers Whose Center Is Not Flat}

Let \(S\) be the spectrum of a discrete valuation ring \(A\), let
\(\pi\) be a uniformizer of \(A\), and let \(t\) and \(s\) be the
generic and closed points of \(S\), respectively. Let \(G\) be the
smooth affine \(S\)-group whose ring is
\[
  B=A[T,T^{-1},U]/F,\qquad F=1-T+\pi U,
\]
the composition law being defined by
\[
  (t,u),(t',u')\longmapsto (tt',\pi uu'+u+u').
\]
The generic fiber \(G_t\) is therefore isomorphic to the
multiplicative group \((\Gm)_t\), while the closed fiber \(G_s\)
is isomorphic to the additive group \((\Ga)_s\). The function
\(T\) is invertible in \(\Gamma(G,\mathcal O_G)\) and defines an
\(S\)-group morphism
\[
  \varphi:G\longrightarrow(\Gm)_S
\]
which is an isomorphism on the generic fiber and the trivial
morphism on the closed fiber. The datum of \(\varphi\) permits one
to construct the semidirect-product group
\[
  H=G\cdot(\Ga)_S
\]
(note that \((\Gm)_S\) acts on \((\Ga)_S\)). The center \(Z\) of
\(H\) is the unit group on the generic fiber and is equal to
\(H_s\) on the closed fiber, and is therefore not flat over \(S\).

% END inputs/XVI/tex/components/03_expose_XVI_section3_flat_center.tex

% BEGIN inputs/XVI/tex/components/04_expose_XVI_section4_affine_fibers.tex
% SGA 3, Exposé XVI, Loop-1 English reconstruction.
% Current scope: all of Section 4.
% Authority: Polo--Gille Expo16.pdf, local pages 18--19
% (combined-reader pages 956--957; printed pages 519--522).

\SGARefTarget{sga3:xvi:section:4}%
\section*{4. Groups with Affine Fibers, of Unipotent Rank Zero}
\addcontentsline{toc}{section}{4. Groups with Affine Fibers,
of Unipotent Rank Zero}

\noindent\SGARefTarget{sga3:xvi:theorem:4:1}%
\textbf{Theorem 4.1.} --- \textit{Let \(S\) be a prescheme and
let \(G\) be an \(S\)-group prescheme smooth over \(S\), of finite
presentation, with affine, connected fibers, of unipotent rank
zero (\SGARefLink{sga3:xv:definition:6:1:ter}{Exposé XV, 6.1 ter}). Then:}

\begin{enumerate}
\renewcommand{\labelenumi}{\textit{\alph{enumi})}}
\item \textit{The center \(Z\) of \(G\) is a group of
multiplicative type (and is therefore a reductive center of \(G\);
\SGARefLink{sga3:xii:4.1}{Exposé XII, 4.1}).}

\item \textit{\(G\) satisfies conditions a) through d) of \SGARefLink{sga3:xvi:theorem:3:3}{3.3}.}

\item \textit{Locally for the étale topology, \(G\) possesses
maximal tori, and these are also Cartan subgroups of \(G\). The
functor \(\mathscr{TM}_G\) of maximal tori of \(G\) is
representable by an \(S\)-prescheme smooth and affine over \(S\).}

\item \textit{\(G\) is quasi-affine over \(S\). Moreover, if
\(T\) is a maximal torus of \(G\), if}
\[
  N=\Norm_G(T)
\]
\textit{is its normalizer in \(G\), and if \(W=N/T\) is the Weyl
group relative to \(T\), the following conditions are equivalent:}
\begin{enumerate}
\renewcommand{\labelenumii}{\textit{(\roman{enumii})}}
\item \textit{\(G\) is affine over \(S\);}
\item \textit{\(G'=G/Z\) is affine over \(S\);}
\item \textit{\(N\) is affine over \(S\);}
\item \textit{\(W\) is affine over \(S\).}
\end{enumerate}
\textit{These conditions are always satisfied if \(S\) is locally
noetherian of dimension at most \(1\).}
\end{enumerate}

\medskip
\noindent\textit{Proof.}
Since \(G\) has affine fibers (and hence abelian rank zero) and
unipotent rank zero, \(G\) possesses maximal tori locally for the
étale topology (\SGARefLink{sga3:xv:corollary:8-18}{Exposé XV, 8.18}), and every maximal torus of \(G\)
is plainly a Cartan subgroup of \(G\). This proves the first part
of c). To see that \(Z\) is of multiplicative type, we may suppose,
by what precedes, that \(G\) possesses a maximal torus \(T\).
Since \(T\) is also a Cartan subgroup, it contains \(Z\) (because
\(T=\Centr_G(T)\) by \SGARefLink{sga3:xii:7.1}{Exposé XII, 7.1}), and
\[
  Z=Z\cap T
\]
is of multiplicative type by \SGARefLink{sga3:xvi:proposition:3:5}{3.5} i). This proves a). Assertion b)
is clear in view of a). On the other hand, the functor
\(\mathscr{TM}_G\) is isomorphic to the functor of Cartan
subgroups of \(G\) (\SGARefLink{sga3:xii:7.1}{Exposé XII, 7.1}), and is therefore smooth and
affine over \(S\) by \SGARefLink{sga3:xvi:theorem:3:3}{3.3} c). This finishes the proof of c). It
remains to prove d).

\medskip
\noindent\textit{Proof of d).}
Since \(G'\) is quasi-affine (\SGARefLink{sga3:xvi:theorem:3:3}{3.3} b)) and \(Z\) is of
multiplicative type, hence affine over \(S\), \(G\) is quasi-affine
(Exposé VI\textsubscript{B}, \S9).

\medskip
\noindent i) \(\Longleftrightarrow\) ii).
If \(G\) is affine, then \(G'\) is affine by \SGARefLink{sga3:ix:proposition:2:3}{Exposé IX, 2.3}. If
\(G'\) is affine, then \(G\) is affine as an extension of an affine
group by an affine group (Exposé VI\textsubscript{B}, \S9).

\medskip
\noindent i) \(\Longleftrightarrow\) iii).
If \(G\) is affine, then \(N\) is affine, since it is closed in
\(G\) (\SGARefLink{sga3:xi:corollary:6:11}{Exposé XI, 6.11} a)). Moreover, \(G/N\) is isomorphic to the
functor \(\mathscr{TM}_G\) (conjugacy of maximal tori; cf.\
\SGARefLink{sga3:xii:7.1}{Exposé XII, 7.1} b)), and is therefore affine over \(S\) by c).
Thus if \(N\) is affine over \(S\), then \(G\) is affine, as an
``extension'' of an affine homogeneous space by an affine group
(Exposé VI\textsubscript{B}, \S9).

\medskip
\noindent iii) \(\Longleftrightarrow\) iv).
If \(N\) is affine, then \(W=N/T\) is affine (\SGARefLink{sga3:ix:proposition:2:3}{Exposé IX, 2.3}), and
the converse follows from Exposé VI\textsubscript{B}, \S9.

Moreover, \(N/T\) is representable by an \(S\)-group prescheme
that is étale, separated, and of finite type over \(S\) (\SGARefLink{sga3:xv:proposition:7-1}{Exposé XV,
7.1} iv)), hence quasi-finite over \(S\). The last assertion of d)
therefore follows from the next lemma, applied with \(X=W\) and
\(Y=S\).

\medskip
\noindent\SGARefTarget{sga3:xvi:lemma:4:2}%
\textbf{Lemma 4.2.} --- \textit{Let \(S\) be a locally noetherian
prescheme of dimension \(1\), let \(X\) and \(Y\) be two
\(S\)-preschemes locally of finite type over \(S\), and let
\(u:X\to Y\) be a quasi-finite and separated \(S\)-morphism.
Suppose that, for every point \(s\) of \(S\), the morphism}
\[
  u_s:X_s\longrightarrow Y_s
\]
\textit{induced from \(u\) by the base change
\(\operatorname{Spec}\kappa(s)\to S\) is finite. Then \(u\) is an
affine morphism.}

\medskip
The assertion to be proved is local on \(Y\), which allows us to
suppose that \(Y\) (and hence also \(X\)) is of finite type over
\(S\). By EGA IV, \S8, it suffices to prove 4.2 when \(S\) is the
spectrum of a local ring \(A\). By fpqc descent we may suppose
that \(A\) is complete, then that \(A\) is reduced (EGA II,
1.6.4), and then that \(A\) is normal (Nagata's theorem, EGA
0\textsubscript{IV}, \S22, and Chevalley's theorem, EGA II,
6.7.1). If \(A\) is a field, then \(u\) is finite by hypothesis,
hence affine. Otherwise \(A\) is a discrete valuation ring; let
\(s\) and \(t\) be the closed and generic points of \(S\), and let
\(\pi\) be a uniformizer of \(A\). Let \(y\) be a point of \(Y_s\).
Applying EGA IV, \S8, once more, we may replace \(Y\) by the
spectrum \(Y'\) of \(\mathcal O_{Y,y}\), and \(X\) by
\[
  X'=X\times_Y Y';
\]
we may even suppose that \(\mathcal O_{Y,y}\) is complete. Since
\(u\) is quasi-finite and separated, EGA II, 6.2.6, shows that
\(X'\) is the disjoint union of two schemes \(X_1\) and \(X_2\),
where \(X_1\) is finite over \(Y'\) (hence affine), while
\(u(X_2)\) does not contain the closed point \(y\) of \(Y'\). I
claim that, under these conditions, \(X_2\) does not meet the
closed fiber \(X'_s\). Indeed, by hypothesis,
\[
  X'_s\longrightarrow Y'_s
\]
is finite, so the restriction of this morphism to the closed
subset \((X_2)_s\) is finite, and its image in \((Y')_s\) is
closed. Since this image does not contain the closed point of the
local scheme \((Y')_s\), \((X_2)_s\) is empty. Since \(u_t\) is
finite, the restriction to \((X_2)_t=X_2\) of the morphism
\[
  X'_t\longrightarrow Y'_t
\]
is finite. Moreover \(Y'_t=Y'_\pi\), so the open immersion
\(Y'_t\to Y'\) is affine. In short, the composite morphism
\[
  X_2\longrightarrow Y'_t\longrightarrow Y'
\]
is affine, and it follows that
\[
  X'=X_1\amalg X_2\longrightarrow Y'
\]
is affine.

\medskip
\noindent\SGARefTarget{sga3:xvi:corollary:4:3}%
\textbf{Corollary 4.3.} --- \textit{Let \(S\) be a prescheme, and
let \(G\) be an \(S\)-group prescheme smooth and of finite
presentation over \(S\), with affine, solvable, connected fibers,
of unipotent rank zero. Then \(G\) is affine over \(S\). If,
moreover, the center of \(G\) is the unit group and \(S\) is
quasi-compact, the canonical morphism}
\[
  i_n:G\longrightarrow\operatorname{GL}_S(G^{(n)})
  \qquad\textit{(cf.\ 3.1 b))}
\]
\textit{is a closed immersion for all sufficiently large \(n\).}

\medskip
To prove that \(G\) is affine over \(S\), we may suppose that
\(G\) possesses a maximal torus \(T\) (\SGARefLink{sga3:xvi:theorem:4:1}{4.1} c)). Since \(G\) has
solvable fibers, the Weyl group relative to \(T\) is the unit
group (BIBLE, 6, Th.\ 1, Cor.\ 3), and condition \SGARefLink{sga3:xvi:theorem:4:1}{4.1} d) iv) is
satisfied. If the center of \(G\) is the unit group, \(i_n\) is a
monomorphism for all sufficiently large \(n\) (\SGARefLink{sga3:xvi:proposition:3:1}{3.1}), and is
therefore a closed immersion (\SGARefLink{sga3:xvi:corollary:1:5}{1.5} b)).

% END inputs/XVI/tex/components/04_expose_XVI_section4_affine_fibers.tex

% BEGIN inputs/XVI/tex/components/05_expose_XVI_section5_reductive_semisimple.tex
% SGA 3, Exposé XVI, Loop-1 English reconstruction.
% Current scope: all of Section 5.
% Authority: Polo--Gille Expo16.pdf, local pages 19--21
% (combined-reader pages 957--959; printed pages 523--525).

\SGARefTarget{sga3:xvi:section:5}%
\section*{5. Application to Reductive and Semisimple Groups}
\addcontentsline{toc}{section}{5. Application to Reductive and
Semisimple Groups}

\noindent\SGARefTarget{sga3:xvi:definition:5:1}%
\textbf{Definition 5.1.} --- \textit{An \(S\)-group prescheme
\(G\) is said to be \emph{reductive} (respectively
\emph{semisimple}) if \(G\) is smooth and affine over \(S\), with
reductive (respectively semisimple) fibers.}

\medskip
\noindent\SGARefTarget{sga3:xvi:5:1:1}%
\textbf{5.1.1.}
Reductive groups will be studied systematically beginning with
Exposé XIX. In this section we shall need the following
properties, which will be proved in Exposé XIX (without using the
developments of the present Exposé).

\begin{enumerate}
\renewcommand{\labelenumi}{\textit{\alph{enumi})}}
\item \textit{Let \(S\) be a prescheme, let \(G\) be an
\(S\)-group prescheme smooth and affine over \(S\), with connected
fibers, and let \(s\) be a point of \(S\) such that \(G_s\) is
reductive (respectively semisimple). Then there exists a
neighborhood \(U\) of \(s\) such that \(G|_U\) is reductive
(respectively semisimple).}

\item \textit{If \(G\) is reductive, then locally for the étale
topology \(G\) possesses maximal tori; and if \(T\) is a maximal
torus of \(G\), the Weyl group}
\[
  W=\Norm_G(T)/T
\]
\textit{is finite over \(S\).}

\item \textit{A reductive group has unipotent rank zero.}
\end{enumerate}

Admitting these facts, we propose to strengthen assertion a)
above.

\medskip
\noindent\SGARefTarget{sga3:xvi:theorem:5:2}%
\textbf{Theorem 5.2.} --- \textit{Let \(S\) be a prescheme, let
\(G\) be an \(S\)-group prescheme smooth and of finite
presentation over \(S\), with connected fibers, and let \(s\) be
a point of \(S\). Then:}

\begin{enumerate}
\renewcommand{\labelenumi}{\textit{(\roman{enumi})}}
\item \textit{If the fibers of \(G\) are affine and \(G_s\) is
reductive, there exists an open neighborhood \(U\) of \(s\) such
that \(G|_U\) is reductive.}

\item \textit{If \(G_s\) is semisimple, there exists an open
neighborhood \(U\) of \(s\) such that \(G|_U\) is semisimple.}
\end{enumerate}

Using \SGARefLink{sga3:xv:proposition:6-2}{Exposé XV, 6.2} i), and Exposé VI\textsubscript{B}, \S10, we
reduce to the case where \(S\) is noetherian.

In case ii), consider the center \(Z\) of \(G\), which is
representable (\SGARefLink{sga3:xi:corollary:6:11}{Exposé XI, 6.11} a)). Since \(G_s\) is semisimple,
it is well known that \(Z_s\) is finite. Consequently (Exposé
VI\textsubscript{B}, \S4), there exists a neighborhood \(U\) of
\(s\) such that \(Z\) is quasi-finite over \(U\). For every point
\(t\) of \(U\), \(G_t\) is then affine (\SGARefLink{sga3:xii:6.1}{Exposé XII, 6.1}).
Shrinking \(S\) if necessary, we may therefore suppose that the
fibers of \(G\) are affine, in case ii) as well as in case i).

In view of a), to prove \SGARefLink{sga3:xvi:theorem:5:2}{5.2} it suffices to prove that \(G\) is
affine over \(S\). Since \(G\) has affine fibers, we know that the
functor of subtori of \(G\) is representable by a smooth
\(S\)-prescheme (\SGARefLink{sga3:xv:corollary:8:11}{Exposé XV, 8.11} and \SGARefLink{sga3:xv:theorem:8:8}{8.9}). After making an étale
extension covering \(s\), which is permissible, we may therefore
suppose that there is a subtorus \(T\) of \(G\) such that \(T_s\)
is a maximal torus of \(G_s\). But then
\[
  C=\Centr_G(T)
\]
is a smooth subgroup of \(G\), with connected fibers, containing
\(T\), such that \(C_s=T_s\) (because \(C_s\) is a Cartan
subgroup of \(G_s\), and \(G_s\) has unipotent rank zero by c)).
It follows that \(C=T\), so \(T\) is a Cartan subgroup and a
maximal torus of \(G\); in particular, the unipotent rank of \(G\)
is zero, and \SGARefLink{sga3:xvi:theorem:4:1}{4.1} applies.

By \SGARefLink{sga3:xvi:theorem:4:1}{4.1} d), it suffices to show that the Weyl group \(W\) relative
to \(T\) is affine over a neighborhood \(U\) of \(s\). In fact,
we shall show that \(W\) is even finite over a neighborhood of
\(s\). Since \(W\) is quasi-finite over \(S\) (\SGARefLink{sga3:xv:proposition:7-1}{Exposé XV,
7.1} iv)), saying that \(W\) is finite over a neighborhood \(U\)
of \(s\) is equivalent to saying that the morphism
\[
  W\longrightarrow S
\]
is proper at \(s\) (EGA IV, 15.7, and EGA III, 4.4.2). To
establish this, the valuative criterion of local properness
(EGA IV, 15.7) reduces us to the case where \(S\) is the spectrum
of a discrete valuation ring, with closed point \(s\). But then,
by the last assertion of \SGARefLink{sga3:xvi:theorem:4:1}{4.1} d), \(G\) is affine over \(S\), and
we may apply property a) recalled above to conclude that \(G\) is
reductive (respectively semisimple). Applying property b), we
conclude that \(W\) is indeed finite over \(S\), which completes
the proof of \SGARefLink{sga3:xvi:theorem:5:2}{5.2}.

% END inputs/XVI/tex/components/05_expose_XVI_section5_reductive_semisimple.tex

% BEGIN inputs/XVI/tex/components/06_expose_XVI_section6_rigidity.tex
% SGA 3, Exposé XVI, Loop-1 English reconstruction.
% Current scope: all of Section 6 and the end of Exposé XVI.
% Authority: Polo--Gille Expo16.pdf, local pages 21--23
% (combined-reader pages 959--961; printed pages 526--530).
% Local authority page 24 (combined-reader page 962) is terminal blank.

\SGARefTarget{sga3:xvi:section:6}%
\section*{6. Applications: Extension of Certain Rigidity Properties
of Tori to Groups of Unipotent Rank Zero}
\addcontentsline{toc}{section}{6. Applications: Extension of Certain
Rigidity Properties of Tori to Groups of Unipotent Rank Zero}

\noindent\SGARefTarget{sga3:xvi:proposition:6:1}%
\textbf{Proposition 6.1.} --- \textit{Let \(S\) be a prescheme and
let \(G\) be an \(S\)-group prescheme smooth over \(S\), of finite
presentation, with connected affine fibers, and whose center is of
multiplicative type (for example, \(G\) of unipotent rank zero;
cf.\ \SGARefLink{sga3:xvi:theorem:4:1}{4.1} a)). Then every closed normal subgroup \(K\) of \(G\),
of finite presentation over \(S\) and quasi-finite over \(S\), is
finite over \(S\).}

\medskip
Indeed, for every geometric point \(\bar s\) above \(S\),
\((K_{\bar s})_{\mathrm{red}}\) is a finite étale subgroup of
\(G_{\bar s}\), normal and therefore central (since \(G_{\bar s}\)
is connected). Consequently
\[
  K'=Z\cap K,
\]
where \(Z\) denotes the center of \(G\), has the same underlying
space as \(K\), and it suffices to show that \(K'\) is finite over
\(S\). Now \(K'\) is a closed quasi-finite subgroup of the group
of multiplicative type \(Z\), and is therefore finite, as is
immediate (locally on \(S\), \(K'\) is contained in \({}_nZ\) for
a suitable integer \(n\), and \({}_nZ\) is finite over \(S\)).

\medskip
\noindent\SGARefTarget{sga3:xvi:corollary:6:2}%
\textbf{Corollary 6.2.} --- \textit{Let \(S\) and \(G\) be as
above, let \(K\) be an \(S\)-group subprescheme of \(G\), of
finite presentation over \(S\), normal and closed in \(G\), and
let \(s\) be a point of \(S\). If \(K_s\) is finite (respectively
is the unit group), there exists an open neighborhood \(U\) of
\(s\) such that \(K|_U\) is finite over \(S\) (respectively is the
unit group).}

\medskip
In view of the upper semicontinuity of the dimension of the fibers
of \(K\) (Exposé VI\textsubscript{B}, \S4), after shrinking \(S\)
we may already suppose that \(K\) is quasi-finite over \(S\), and
hence finite by \SGARefLink{sga3:xvi:proposition:6:1}{6.1}. If, moreover, \(K_s\) is the unit group, then
Nakayama's lemma easily implies that \(K\) is the unit group over
\(\operatorname{Spec}\mathcal O_s\), and hence over a neighborhood
of \(s\).

\medskip
\noindent\SGARefTarget{sga3:xvi:proposition:6:3}%
\textbf{Proposition 6.3.} --- \textit{Let \(S\) be a prescheme,
let \(G\) be an \(S\)-group prescheme smooth and of finite
presentation over \(S\), with affine, connected fibers, of
unipotent rank zero, let \(H\) be an \(S\)-group prescheme, let
\(s\) be a point of \(S\), and let}
\[
  u:G\longrightarrow H
\]
\textit{be an \(S\)-group morphism such that \(\Ker u_s\) is
central. Suppose, moreover, that \(H\) is of finite presentation
over \(S\), or that \(S\) is locally noetherian and \(G\) is
locally of finite type over \(S\). Then there exists an open
neighborhood \(U\) of \(s\) such that, if \(K=\Ker u\), \(K|_U\)
is central and of multiplicative type. Moreover,}
\[
  G'=(G|_U)/(K|_U)
\]
\textit{is representable, and the morphism}
\[
  u':G'\longrightarrow H|_U
\]
\textit{induced from \(u\) by passage to the quotient is an
immersion.}

\medskip
\noindent\textit{Proof.}
Let \(Z\) be the center of \(G\), and put \(K'=K\cap Z\).

\medskip
\noindent a) \(K'\) is a subgroup of multiplicative type over a
neighborhood \(U\) of \(s\). Indeed, after making an étale
extension over a neighborhood \(U\) of \(s\), which is legitimate,
we may suppose that \(G\) possesses a maximal torus \(T\)
(\SGARefLink{sga3:xvi:theorem:4:1}{4.1} c)). But then \(T\cap K\) is a group of multiplicative type
(\SGARefLink{sga3:xv:lemma:8:3}{Exposé XV, 8.3}) whose fiber at \(s\) is central by hypothesis;
therefore \(T\cap K\) is central over a neighborhood of \(s\)
(\SGARefLink{sga3:ix:theorem:5:6}{Exposé IX, 5.6} a)), and consequently coincides with \(K'\).

\medskip
\noindent b) Let us show that \(K'=K\) over a neighborhood \(U\)
of \(s\). By a), we may already suppose that \(K'\) is of
multiplicative type, and hence flat over \(S\). Since \(K'_s=K_s\),
the natural immersion
\[
  K'\longrightarrow K
\]
is then open over a neighborhood \(U\) of \(s\) (Exposé
VI\textsubscript{B}, \SGARefLink{cor:sga3-VIB-2.6}{2.6}). If \(t\in U\), the image of \(K_t\) in
the group
\[
  G'_t=G_t/Z_t
\]
is a finite étale group, normal in \(G'_t\) (since \(K_t\) is
normal in \(G_t\)), hence central, and therefore reduced to the
unit group (\SGARefLink{sga3:xvi:theorem:3:3}{3.3} b)). This says that \(K_t\) is contained in
\(Z_t\), and hence is equal to \(K'_t\), whence \(K=K'\) over
\(U\).

\medskip
\noindent c) The representability of \(G'\) now follows from \SGARefLink{sga3:xvi:corollary:2:3}{2.3};
the fact that \(u'\) is an immersion is contained in \SGARefLink{sga3:xvi:theorem:1:3}{1.3} a), in
view of \SGARefLink{sga3:xvi:theorem:4:1}{4.1} c) applied to \(G'\).

\medskip
\noindent\SGARefTarget{sga3:xvi:proposition:6:4}%
\textbf{Proposition 6.4.} --- \textit{Let \(S\) be a prescheme,
let \(G\) be an \(S\)-group prescheme smooth and of finite
presentation over \(S\), with connected fibers, and let \(s\) be
a point of \(S\) such that \(G_s\) is generated by its subtori
(\SGARefLink{sga3:xii:8.2}{Exposé XII, 8.2}) (for example, \(G_s\) affine of unipotent rank
zero). Let \(H\) be an \(S\)-group prescheme, and let}
\[
  u,v:G\longrightarrow H
\]
\textit{be two \(S\)-homomorphisms such that \(u_s=v_s\) and
\(u_s\) is central. Suppose, moreover, that \(H\) is of finite
presentation over \(S\), or that \(S\) is locally noetherian.
Then there exists a neighborhood \(U\) of \(s\) such that}
\[
  u|_U=v|_U.
\]

\medskip
We reduce as usual to the case where \(S\) is noetherian (to study
the condition ``\(G_s\) is generated by its subtori,'' one uses
Exposé VI\textsubscript{B}, \SGARefLink{prop:sga3-VIB-7.4}{7.4}). Since \(G\) has connected fibers,
to prove that \(u=v\) over a neighborhood \(U\) of \(s\), it
suffices to prove that \(u=v\) after reduction modulo every power
of the maximal ideal of the local ring \(\mathcal O_{S,s}\), which
reduces us to the case where \(S\) is local artinian.

But then the functor of maximal subtori of \(G\) is representable
by an \(S\)-scheme \(X\), smooth and of finite type over \(S\)
(\SGARefLink{sga3:xv:corollary:8:17}{Exposé XV, 8.17}). Let \(T\) be the maximal subtorus of \(G_X\)
universal for the functor \(X\). The hypothesis on \(G_s\) means
that the algebraic subgroup of \(G_s\) generated (Exposé
VI\textsubscript{B}, \S7) by the \(\kappa(s)\)-morphism
\[
  f:T_s\longrightarrow G_s
\]
(the composite of the immersion \(T_s\to G_{X_s}\) and the
canonical projection \(G_{X_s}\to G_s\)) is all of \(G_s\). But
\(X_s\) is geometrically connected (for if \(\overline T\) is a
maximal torus of \(G_{\bar s}\), and \(N\) its normalizer in
\(G_{\bar s}\), then \(X_{\bar s}\) is isomorphic to
\(G_{\bar s}/N\)), and the image of \(T_s\) under \(f\) contains
the unit section of \(G_s\). It follows from Exposé
VI\textsubscript{B}, \SGARefLink{prop:sga3-VIB-7.4}{7.4}, that there exist an integer \(N>0\) and
an \(S\)-morphism
\[
  f^N:T^N\longrightarrow G,
  \qquad
  T^N=T\times_S\cdots\times_S T
  \quad\text{(\(N\) factors),}
\]
where \(f^N\) depends only on \(f\) and the multiplication law of
\(G\), such that \((f^N)_s\) is surjective.

Consider the base change \(X\to S\) and the restrictions of
\(u_X\) and \(v_X\) to the subtorus \(T\) of \(G_X\). By
hypothesis,
\[
  u_{X_s}=v_{X_s},
\]
and
\[
  u_{X_s}:T_{X_s}\longrightarrow H_{X_s}
\]
is a central homomorphism. It then follows from \SGARefLink{sga3:ix:theorem:5:1}{Exposé IX, 5.1},
that
\[
  u_X|_T=v_X|_T.
\]
Making the definition of \(f^N\) explicit, and taking account of
the fact that \(u\) and \(v\) are homomorphisms, one immediately
deduces that
\[
  uf^N=vf^N.
\]
But the \(\kappa(s)\)-morphism \(f_s^N\) is surjective and \(G_s\)
is smooth, hence reduced; consequently \(f_s^N\) is generically
flat. Since \(T\) is smooth over \(S\), hence flat, there exists a
nonempty open subset \(V\) of \(T^N\) such that
\[
  f^N|_V
\]
is flat (EGA IV, 11.3.10).\footnote{The French source prints \(T\)
here; \(T^N\) is the domain of \(f^N\).} The image of \(V\) is an open subset
\(W\) of \(G\) (EGA IV, 11.3.1), and
\[
  f^N:V\longrightarrow W
\]
is faithfully flat and of finite presentation, hence covering for
the fpqc topology. The equality \(uf^N=vf^N\) then implies
\[
  u|_W=v|_W.
\]
Since \(W\) is schematically dense in \(G\) (EGA IV, 11.10.10),
and \(H\) is separated over \(S\) (Exposé VI\textsubscript{A},
\SGARefLink{sga3:VIA:0.3}{0.3}), it follows immediately that \(u=v\).

% END inputs/XVI/tex/components/06_expose_XVI_section6_rigidity.tex


